%%%%%%%%%%%%%%%%%%%%%%%%%%%%%%%%%%%%%%%%%%%%%%%%%%%%%%%%%%%%%%%%%%%%%%
% Lớp 11 | Giới hạn. Hàm số liên tục | Giới hạn của hàm số
% Dạng 1D3?2-5: PP lượng liên hợp, kết quả hữu hạn
% Số câu: 143  (sinh tự động bởi tach_cau.py; từ nay thầy sửa trực tiếp file này)
%%%%%%%%%%%%%%%%%%%%%%%%%%%%%%%%%%%%%%%%%%%%%%%%%%%%%%%%%%%%%%%%%%%%%%

%%%%%%%%%%%%%%%% TRẮC NGHIỆM 4 PHƯƠNG ÁN %%%%%%%%%%%%%%%%

% nguồn: Ok_Toan11_TracNghiem_Lan2
\begin{ex}%[1D3H2-5]%[ma:03f3a783]
	Biết  $\lim\limits_{x\to -\infty } \left( \sqrt{5x^2+2x}+x\sqrt{5} \right)=a\sqrt{5}+b$ với $a$, $b\in \mathbb{Q}$. Tính $S=5a+b$.
	 \choice
		{$S=-5$} 
		 {$S=-1$}
		{\True $S=1$} 
		{$S=5$} 
	\loigiai{ 
		Ta có $\lim\limits_{x\to -\infty } \left( \sqrt{5x^2+2x}+x\sqrt{5} \right)=\lim\limits_{x\to -\infty }  \dfrac{2x}{\sqrt{5x^2+2x}-x\sqrt{5}}=\lim\limits_{x\to -\infty } \dfrac{2}{-\sqrt{5+\dfrac{2}{x}}-\sqrt{5}}=-\dfrac{1}{5}\sqrt{5}$.\\
	Suy ra $a=-\dfrac{1}{5}$, $b=0$. Vậy $S=-1$.} 
\end{ex}

% nguồn: Ok_Toan11_TracNghiem_Lan2
\begin{ex}%[1D3H2-5]%[ma:02461d85]
	 Cho $\lim\limits_{x\to -\infty } \left( \sqrt{9x^2+ax}+3x \right)=-2$. Tính giá trị của $a$.
	\choice 
		 {$-6$}
		{\True $12$} 
		 {$6$}
		 {$-12$}
	 \loigiai{
		Ta có $\lim\limits_{x\to -\infty } \left( \sqrt{9x^2+ax}+3x \right)=\lim\limits_{x\to -\infty } \left( \dfrac{ax}{\sqrt{9x^2+ax}-3x}  \right)=\lim\limits_{x\to -\infty } \dfrac{a}{-\sqrt{9+\dfrac{a}{x}}-3}=-\dfrac{a}{6}$.\\
		Suy ra $-\dfrac{a}{6}=-2\Leftrightarrow  a=12$.
	 }
\end{ex}

% nguồn: Ok_Toan11_TracNghiem_Lan2
\begin{ex}%[1D3H2-5]%[ma:15cd1cb7]
	 Cho $I=\lim\limits_{x\to 0} \dfrac{2\left(\sqrt{3x+1}-1\right)}{x}$ và $J=\lim\limits_{x\to-1} \dfrac{x^2-x-2}{x+1}$. Tính $I-J$.
	 \choice
		{\True $6$} 
		 {$3$}
		{$-6$} 
		 {$0$}
	 \loigiai{
		Ta có 
		\begin{itemize}
			\item $I=\lim\limits_{x\to 0}  \dfrac{2\left(\sqrt{3x+1}-1\right)}{x}=\lim\limits_{x\to 0} \dfrac{6x}{x\left(\sqrt{3x+1}+1\right)}=\lim\limits_{x\to 0} \dfrac{6}{\sqrt{3x+1}+1}=3$.
			\item $J=\lim\limits_{x\to-1}  \dfrac{x^2-x-2}{x+1}=\lim\limits_{x\to-1} \dfrac{\left(x+1\right)\left(x-2\right)}{x+1}=\lim\limits_{x\to-1} \left(x-2\right)=-3$.
		\end{itemize}
		 Khi đó $I-J=6$.
	} 
\end{ex}

% nguồn: Ok_Toan11_TracNghiem_Lan2
\begin{ex}%[1D3H2-5]%[ma:6adbc6fd]
	 Tính $\lim\limits_{x\to 1}\dfrac{\sqrt{5-x}-2}{x^2-1}$.
	 \choice
		 {$\dfrac{1}{8}$}
		 {$\dfrac{1}{4}$}
		{\True $-\dfrac{1}{8}$} 
		{$-\dfrac{1}{4}$} 
	\loigiai{Ta  có
		$\lim\limits_{x\to 1}\dfrac{\sqrt{5-x}-2}{x^2-1}=\lim\limits_{x\to 1}\dfrac{5-x-4}{(x-1)(x+1)(\sqrt{5-x}+2)}=\lim\limits_{x\to  1}\dfrac{-(x-1)}{(x-1)(x+1)(\sqrt{5-x}+2)}=-\dfrac{1}{8}$.
	 }
\end{ex}

% nguồn: Ok_Toan11_TracNghiem_Lan2
\begin{ex}%[1D3H2-5]%[ma:c2fae5c4]
	 Biết $\lim\limits_{x\to 0}\dfrac{\sqrt{x+1}-1}{\sqrt{2x+1}-1}=\dfrac{a}{b}$ ( với $a$, $b$ nguyên dương và nguyên tố cùng nhau). Khẳng định nào sao đây là đúng?
	 \choice
		{\True  $a+b=3$}
		{$a+b=-3$} 
		{$a+b=2$} 
		{$a+b=1$} 
	 \loigiai{Ta có
		\begin{align*}
			 &\lim\limits_{x\to 0}\dfrac{\sqrt{x+1}-1}{\sqrt{2x+1}-1}\\
			 =&\lim\limits_{x\to 0}\dfrac{(\sqrt{x+1}-1)(\sqrt{x+1}+1)(\sqrt{2x+1}+1)}{(\sqrt{2x+1}-1)(\sqrt{2x+1}+1)(\sqrt{x+1}+1)}\\
			=&\lim\limits_{x\to  0}\dfrac{x(\sqrt{2x+1}+1)}{2x(\sqrt{x+1}+1)}\\
			 =&\lim\limits_{x\to 0}\dfrac{\sqrt{2x+1}+1}{2(\sqrt{x+1}+1)}\\
			 =&\dfrac{1}{2}.
		\end{align*}
		Suy ra $a=1$; $b=2$.\\ 
		Vậy $a+b=3$. 
	 }
\end{ex}

% nguồn: Ok_Toan11_TracNghiem_Lan2
\begin{ex}%[1D3H2-5]%[ma:64b64486]
	Tính $M=\lim\limits_{x\to-\infty} \left(\sqrt{x^2-4x}-\sqrt{x^2-x} \right)$. 
	 \choice
		{$-\dfrac{3}{2}$} 
		 {$\dfrac{1}{2}$}
		 {\True $\dfrac{3}{2}$}
		{$-\dfrac{1}{2}$} 
	 \loigiai{
		Ta có \begin{align*}
		M &= \lim\limits_{x\to-\infty} \left(\sqrt{x^2-4x}-\sqrt{x^2-x} \right)  \\
		&= \lim\limits_{x\to-\infty} \dfrac{-3x}{\sqrt{x^2-4x}+\sqrt{x^2-x}} \\ 
		 &= \lim\limits_{x\to-\infty} \dfrac{-3x}{\left|x \right|\cdot\left(\sqrt{1-\dfrac{4}{x}}+\sqrt{1-\dfrac{1}{x}} \right)} \\
		 &= \lim\limits_{x\to-\infty} \dfrac{3}{\sqrt{1-\dfrac{4}{x}}+\sqrt{1-\dfrac{1}{x}}} \\
		&=  \dfrac{3}{2}.
	\end{align*}
 }
\end{ex}

% nguồn: Ok_Toan11_TracNghiem_Lan2
\begin{ex}%[1D3H2-5]%[ma:78fa03f7]
	Kết quả của giới hạn $\lim\limits_{x\to-\infty}\left(\sqrt{9x^2+8x+2020}+3x\right)$  là
	\choice 
		{$-\infty$} 
		{$\dfrac{4}{3}$} 
		 {\True $\dfrac{-4}{3}$}
		 {$+\infty$}
	\loigiai{ 
	$\begin{aligned}
	\lim\limits_{x\to-\infty}\left(\sqrt{9 x^2+8 x+2020}+3 x\right)&=\lim\limits_{x\to-\infty} \dfrac{9 x^2+8 x+2020-9 x^2}{\sqrt{9  x^2+8 x+2020}-3 x}\\&=\lim\limits_{x\to-\infty} \dfrac{8 x+2020}{|x|\sqrt{9+\dfrac{8}{x}+\dfrac{2020}{x^2}}-3 x}\\
	&=\lim\limits_{x\to-\infty} \dfrac{8 x+2020}{-x \sqrt{9+\dfrac{8}{x}+\dfrac{2020}{x^2}}-3 x}\\&=\lim\limits_{x\to-\infty}  \dfrac{8+\dfrac{2020}{x}}{-\sqrt{9+\dfrac{8}{x}+\dfrac{2020}{x^2}}-3}=-\dfrac{4}{3}.
\end{aligned}$\\
 Vậy $\lim\limits_{x\to-\infty}\left(\sqrt{9 x^2+8 x+2020}+3 x\right)=-\dfrac{4}{3}$.
 }
\end{ex}

% nguồn: Ok_Toan11_TracNghiem_Lan2
\begin{ex}%[1D3H2-5]%[ma:6aec0ead]
	$\lim\limits_{x\to +\infty }  x\left( \sqrt{x^2+5x+4}-\sqrt{x^2+5x-2} \right)$ bằng
	\choice 
		{\True $3$} 
		{$1$} 
		{$0$} 
		{$+\infty  $}
	 \loigiai{
		Ta  có
		\begin{align*}
			 &\lim\limits_{x\to +\infty } x\left( \sqrt{x^2+5x+4}-\sqrt{x^2+5x-2} \right)\\
			&=\lim\limits_{x\to +\infty}  \dfrac{6x}{\sqrt{x^2+5x+4}+\sqrt{x^2+5x-2}}\\
			 &=\lim\limits_{x\to +\infty } \dfrac{6x}{x\left( \sqrt{1+\dfrac{5}{x}+\dfrac{4}{x^2}}+\sqrt{1+\dfrac{5}{x}-\dfrac{2}{x^2}} \right)}\\
			&=3. 
	\end{align*}}
\end{ex}

% nguồn: Ok_Toan11_TracNghiem_Lan2
\begin{ex}%[1D3H2-5]%[ma:cb179784]
	$\lim\limits_{x\to -\infty } \left(  \sqrt{4x^2+8x+1}+2x \right)$ bằng
	\choice 
		{$-\infty  $}
		 {$0$}
		{\True  $-2$}
		 {$+\infty $}
	 \loigiai{
		Ta có 
		\[ 
		\lim\limits_{x\to  -\infty } (\sqrt{4x^2+8x+1}+2x)=\lim\limits_{x\to -\infty } \dfrac{8x+1}{\sqrt{4x^2+8x+1}-2x}=\lim\limits_{x\to -\infty } \dfrac{8+\dfrac{1}{x}}{-\sqrt{4+\dfrac{8}{x}+\dfrac{1}{x^2}}-2}=-2.
	 \]}
\end{ex}

% nguồn: Ok_Toan11_TracNghiem_Lan2
\begin{ex}%[1D3H2-5]%[ma:fc311982]
	Cho hàm số $f(x)=\dfrac{\sqrt{4x+1}-3}{x-2}$. Tính $\lim\limits_{x\to 2} f(x)$. 
	 \choice
		 {$\lim\limits_{x\to 2} f(x)=-\dfrac{2}{3}$}
		{$\lim\limits_{x\to  2} f(x)=-\dfrac{3}{2}$}
		{\True $\lim\limits_{x\to 2} f(x)=\dfrac{2}{3}$} 
		{$\lim\limits_{x\to  2} f(x)=\dfrac{3}{2}$}
	 \loigiai{
		Ta  có
		\begin{align*}
			 \lim\limits_{x\to 2} f(x) &=\lim\limits_{x\to 2} \dfrac{\sqrt{4x+1}-3}{x-2}\\
			&=\lim\limits_{x\to 2}  \dfrac{4\left(x-2\right)}{\left(x-2\right)\left(\sqrt{4x+1}+3\right)}\\
			&=\lim\limits_{x\to 2} \dfrac{4}{\sqrt{4x+1}+3}=\dfrac{2}{3}. 
		\end{align*}
	 }
\end{ex}

% nguồn: Ok_Toan11_TracNghiem_Lan2
\begin{ex}%[1D3H2-5]%[ma:f65555ee]
	Tìm giới hạn $I=\lim\limits_{x\to-\infty}\left(\sqrt{x^2+4x+1}+x\right)$. 
	 \choice
		{\True  $I=-2$}
		 {$I=-4$}
		 {$I=1$}
		{$I=-1$} 
	\loigiai{ 
	\[\begin{aligned}
	I&=\lim\limits_{x\to-\infty}\left(\sqrt{x^2+4 x+1}+x\right)=\lim\limits_{x\to-\infty}\left[\dfrac{\left(\sqrt{x^2+4  x+1}+x\right)\left(\sqrt{x^2+4 x+1}-x\right)}{\sqrt{x^2+4 x+1}-x}\right]\\
	&=\lim\limits_{x\to-\infty}\left(\dfrac{4+\dfrac{1}{x}}{\dfrac{-x \sqrt{1+\dfrac{4}{x}+\dfrac{1}{x^2}}-x}{x}}\right)=\dfrac{4}{-2}=-2. 
\end{aligned}\]
 }
\end{ex}

% nguồn: Ok_Toan11_TracNghiem_Lan2
\begin{ex}%[1D3H2-5]%[ma:6bf2acce]
	Tìm $\lim\limits_{x\to -\infty } \left(  \sqrt{x^2+x+2}+x+2 \right)$.
	\choice 
		{\True $\dfrac{3}{2}$} 
		 {$0$}
		 {$-\infty $}
		{$-2$} 
	\loigiai{ 
		Ta có \begin{align*}
		&\lim\limits_{x\to -\infty } \left( \sqrt{x^2+x+2}+x+2 \right)  \\
		&= \lim\limits_{x\to  -\infty } \dfrac{x^2+x+2-{{\left( x+2 \right)}^2}}{\sqrt{x^2+x+2}-x-2} \\
		&= \lim\limits_{x\to -\infty } \dfrac{-3x-2}{\sqrt{x^2+x+2}-x-2}  \\
		&= \lim\limits_{x\to -\infty } \dfrac{-3-\dfrac{2}{x}}{-\sqrt{1+\dfrac{1}{x}+\dfrac{2}{x^2}}-1-\dfrac{2}{x}} \\ 
		&= \dfrac{3}{2}. 
	\end{align*}
} 
\end{ex}

% nguồn: Ok_Toan11_TracNghiem_Lan2
\begin{ex}%[1D3H2-5]%[ma:0eec2dee]
	 Tìm $\lim\limits_{x\to 1}\dfrac{x-\sqrt{2x-1}}{x^2+x-2}$.
	\choice 
		{$-\dfrac{1}{3}$} 
		{$-\infty$} 
		 {\True $0$}
		{$1$} 
	 \loigiai{
		Ta  có
		 $\lim\limits_{x\to 1}\dfrac{x-\sqrt{2x-1}}{x^2+x-2}=\lim\limits_{x\to 1}\dfrac{x^2-2x+1}{(x-1)(x+2)(x+\sqrt{2x-1})}=\lim\limits_{x\to 1}\dfrac{x-1}{(x+2)(x+\sqrt{2x-1})}=0$.
	} 
\end{ex}

% nguồn: Ok_Toan11_TracNghiem_Lan2
\begin{ex}%[1D3H2-5]%[ma:0b943525]
	Tìm $\lim\limits_{x\to 2}\dfrac{x^2-5x+6}{\sqrt{4x+1}-3}$. 
	\choice 
		{$\dfrac{3}{2}$} 
		 {$-\dfrac{1}{2}$}
		{\True  $-\dfrac{3}{2}$}
		{$\dfrac{1}{2}$} 
	\loigiai{ Ta có 
		$\lim\limits_{x\to 2}\dfrac{x^2-5x+6}{\sqrt{4x+1}-3}=\lim\limits_{x\to  2}\dfrac{(x-2)(x-3)(\sqrt{4x+1}+3)}{4(x-2)}=\lim\limits_{x\to 2}\dfrac{(x-3)(\sqrt{4x+1}+3)}{4}=-\dfrac{3}{2}$.
	 }
\end{ex}

% nguồn: Ok_Toan11_TracNghiem_Lan2
\begin{ex}%[1D3H2-5]%[ma:72f90fd1]
	Tìm $\lim\limits_{x\to -\infty  } \left( \sqrt{x^2+x}+2x \right)$.
	\choice 
		 {$2$}
		{\True $-\infty  $}
		 {$1$}
		{$+\infty $} 
	 \loigiai{
		Ta có \begin{align*}
		&\lim\limits_{x\to -\infty }  \left( \sqrt{x^2+x}+2x \right) \\
		&= \lim\limits_{x\to -\infty } \left( \left| x  \right|\sqrt{1+\dfrac{1}{x}}+2x \right) \\
		&= \lim\limits_{x\to -\infty } \left( -x\sqrt{1+\dfrac{1}{x}}+2x \right)  \\
		&= \lim\limits_{x\to -\infty  } \left( x\left( 2-\sqrt{1+\dfrac{1}{x}} \right) \right) \\
		 &= -\infty.
	\end{align*}
vì $\lim\limits_{x\to  -\infty } x=-\infty $ và $\lim\limits_{x\to -\infty } \left( 2-\sqrt{1+\dfrac{1}{x}} \right)=1$.}
\end{ex}

% nguồn: Ok_Toan11_TracNghiem_Lan2
\begin{ex}%[1D3H2-5]%[ma:89acc023]
	Tìm $I=\lim\limits_{x\to  -\infty } \left( \sqrt{x^2+4x+1}+x \right)$.
	 \choice
		 {\True $I=-2$}
		{$I=-4$} 
		 {$I=1$}
		 {$I=-1$}
	\loigiai{ 
		Cách 1: Sử dụng máy tính cầm tay tính giá trị  biểu thức $\sqrt{x^2+4x+1}+x$ tại $x=-{{10}^{10}}$.
		\\ 
		Vậy $I=\lim\limits_{x\to  -\infty } \left( \sqrt{x^2+4x+1}+x \right)$ $=-2$.
		 \\
		Cách 2:  Ta có
		\begin{align*}
			I  &= \lim\limits_{x\to -\infty } \left( \sqrt{x^2+4x+1}+x \right) \\
			&= \lim\limits_{x\to -\infty } \dfrac{4x+1}{\sqrt{x^2+4x+1}-x} \\ 
			&= \lim\limits_{x\to  -\infty } \dfrac{4+\dfrac{1}{x}}{-\sqrt{1+\dfrac{4}{x}+\dfrac{1}{x^2}}-1} \\
			&=  \dfrac{4}{-2} \\
			 &= -2.
	\end{align*}}
\end{ex}

% nguồn: Ok_Toan11_TracNghiem_Lan2
\begin{ex}%[1D3H2-5]%[ma:0c999c5c]
	 Cho hàm số $f(x)=\dfrac{\sqrt{x^2+3a^2}-2a}{x-a}+a$. Tính $\lim\limits_{x\to a}f(x)$.
	 \choice
		{$\lim\limits_{x\to  a}f(x)=\dfrac{2a-1}{2}$}
		 {\True $\lim\limits_{x\to a}f(x)=\dfrac{2a+1}{2}$}
		{$\lim\limits_{x\to a}f(x)=\dfrac{2}{2a+1}$} 
		 {$\lim\limits_{x\to a}f(x)=\dfrac{2}{2a-1}$}
	 \loigiai{
		 Ta có
		\begin{align*}
			\lim\limits_{x\to  a}f(x)&=\lim\limits_{x\to a}\left(\dfrac{\sqrt{x^2+3a^2}-2a}{x-a}+a\right)\\
			 &=\lim\limits_{x\to a}\left(\dfrac{x^2-a^2}{(x-a)(\sqrt{x^2+3a^2}+2a)}+a\right)\\
			 &=\lim\limits_{x\to a}\left(\dfrac{x+a}{\sqrt{x^2+3a^2}+2a}+a\right)=\dfrac{2a}{4a}+a\\
			 &=\dfrac{1}{2}+a\\
			&=\dfrac{2a+1}{2}. 
		\end{align*}
	} 
\end{ex}

% nguồn: Ok_Toan11_TracNghiem_Lan2
\begin{ex}%[1D3H2-5]%[ma:c3f7a760]
	Tính  $\lim\limits_{x\to 0}\dfrac{1-\sqrt{x+1}}{3x}$ bằng
	\choice 
		 {$-\dfrac{1}{3}$}
		 {$\dfrac{1}{3}$}
		 {$\dfrac{1}{6}$}
		{\True  $-\dfrac{1}{6}$}
	\loigiai{Ta  có
		$\lim\limits_{x\to 0}\dfrac{1-\sqrt{x+1}}{3x}=\lim\limits_{x\to 0}\dfrac{1-(x+1)}{3x(1+\sqrt{x+1})}=\lim\limits_{x\to  0}\dfrac{-x}{3x(1+\sqrt{x+1})}=-\dfrac{1}{3\cdot \left(1+\sqrt{0+1}\right)}=-\dfrac{1}{6}$.
	} 
\end{ex}

% nguồn: Ok_Toan11_TracNghiem_Lan2
\begin{ex}%[1D3H2-5]%[ma:2ba25592]
	Biết rằng $\lim\limits_{x\to -\infty  } \left( \sqrt{2x^2-3x+1}+x\sqrt{2} \right)=\dfrac{a}{b}\sqrt{2}$, ($a$, $b\in \mathbb{Z}$, $\dfrac{a}{b}$ tối giản). Tổng $a+b$ có giá trị là
	\choice 
		{$1$} 
		{$5$} 
		{$4$} 
		{\True  $7$}
	\loigiai{ 
		Ta có \begin{align*}
		&\lim\limits_{x\to -\infty }  \left( \sqrt{2x^2-3x+1}+x\sqrt{2} \right) \\
		&=  \lim\limits_{x\to -\infty } \dfrac{2x^2-3x+1-2x^2}{\sqrt{2x^2-3x+1}-x\sqrt{2}} \\
		&= \lim\limits_{x\to -\infty  } \dfrac{x\left( -3+\dfrac{1}{x} \right)}{x\left( -\sqrt{2-\dfrac{3}{x}+\dfrac{1}{x^2}}-\sqrt{2} \right)} \\
		&= \lim\limits_{x\to  -\infty } \dfrac{-3+\dfrac{1}{x}}{-\sqrt{2-\dfrac{3}{x}+\dfrac{1}{x^2}}-\sqrt{2}} \\
		&= \dfrac{3\sqrt{2}}{4}. 
	\end{align*}
Vậy  $a=3$; $b=4\Rightarrow a+b=7$.}
\end{ex}

% nguồn: Ok_Toan11_TracNghiem_Lan2
\begin{ex}%[1D3H2-5]%[ma:a46853d3]
	Biết $\lim\limits_{x\to -\infty } \left( \sqrt{4x^2+ax+1}+bx \right)=-1$.  Tính giá của biểu thức $P=a^2-2b^3$.
	 \choice
		 {$P=32$}
		 {$P=0$}
		{\True  $P=16$}
		 {$P=8$}
	 \loigiai{
		TH1: Với \begin{align*}
		b=2 &\Rightarrow \lim\limits_{x\to -\infty } \left( \sqrt{4x^2+ax+1}+2x  \right) \\
		 &= \lim\limits_{x\to -\infty } \dfrac{ax+1}{\sqrt{4x^2+ax+1}-2x} \\
		&= \lim\limits_{x\to -\infty } \dfrac{a+\dfrac{1}{x}}{-\sqrt{4+\dfrac{a}{x}+\dfrac{1}{x^2}}-2} \\ 
		&=  -\dfrac{a}{4}.
	\end{align*}
	Suy ra $\lim\limits_{x\to -\infty }  \left( \sqrt{4x^2+ax+1}+bx \right)=-1\Leftrightarrow -\dfrac{a}{4}=-1\Rightarrow a=4$.\\
	 TH2: Với $b\ne 2$ Suy ra
	\begin{align*}
		\lim\limits_{x\to  -\infty } \left( \sqrt{4x^2+ax+1}+bx \right)
		&= \lim\limits_{x\to -\infty }  \left[x\left( -\sqrt{4+\dfrac{a}{x}+\dfrac{1}{x^2}}+b \right) \right] \\
		&= \heva{ 
		&-\infty \text{nếu } b \le 2 \\ 
		&+\infty \text{nếu } b  \ge 2.}
	\end{align*}
Vậy $a=4$, $b=2\Rightarrow  P=a^2-2b^3=0$.}
\end{ex}

% nguồn: Ok_Toan11_TracNghiem_Lan2
\begin{ex}%[1D3H2-5]%[ma:9a9a5057]
	$\lim\limits_{x\to  1}\dfrac{\sqrt{x+3}-2}{\sqrt{2x-1}-1}=\dfrac{a}{b}$ ( với $a$, $b$ nguyên dương và nguyên tố cùng nhau). Khẳng định nào sau đây là đúng?
	\choice 
		{$a+b=-5$} 
		{$a+b=2$} 
		 {$a+b=1$}
		{\True $a+b=5$} 
	 \loigiai{
		 Ta có
		\begin{align*}
			&\lim\limits_{x\to 1}\dfrac{\sqrt{x+3}-2}{\sqrt{2x-1}-1}\\ 
			=&\lim\limits_{x\to  1}\dfrac{(\sqrt{x+3}-2)(\sqrt{x+3}+2)(\sqrt{2x-1}+1)}{(\sqrt{2x-1}-1)(\sqrt{2x-1}+1)(\sqrt{x+3}+2)}\\
			=&\lim\limits_{x\to  1}\dfrac{(x-1)(\sqrt{2x-1}+1)}{(2x-2)(\sqrt{x+3}+2)}\\
			 =&\lim\limits_{x\to 1}\dfrac{\sqrt{2x-1}+1}{2(\sqrt{x+3}+2)}\\
			=&\dfrac{1}{4}. 
		\end{align*}
		Do  đó $a=1$; $b=4$.\\
		 Vậy $a+b=5$.
	} 
\end{ex}

% nguồn: Ok_Toan11_TracNghiem_Lan2
\begin{ex}%[1D3H2-5]%[ma:dae4080c]
	Tìm ${\lim\limits_{x\to +\infty } \left( x+1-\sqrt[3]{x^3+2}  \right)}$.
	 \choice
		{$-1$} 
		 {$-\infty $}
		{$+\infty $} 
		 {\True $1$}
	\loigiai{ 
		Ta có \begin{align*}
		\lim\limits_{x\to +\infty } \left( 1+x-\sqrt[3]{x^3+2}  \right)
		&= \lim\limits_{x\to  +\infty } \left( 1+\dfrac{-2}{x^2+x\sqrt[3]{x^3+2}+{{\left( \sqrt[3]{x^3+2} \right)}^2}} \right) \\
		&= \lim\limits_{x\to +\infty } \left( 1+\dfrac{-2}{x^2\left( 1+\sqrt[3]{1+\dfrac{2}{x^3}}+{{\left( \sqrt[3]{1+\dfrac{2}{x^3}} \right)}^2} \right)} \right)  \\
		&= \lim\limits_{x\to +\infty } \left( 1+\dfrac{\dfrac{-2}{x^2}}{1+\sqrt[3]{1+\dfrac{2}{x^3}}+{{\left( \sqrt[3]{1+\dfrac{2}{x^3}} \right)}^2}} \right)  \\
		 &= 1.
	\end{align*}
 Vậy $\lim\limits_{x\to +\infty } \left( x+1-\sqrt[3]{x^3+2} \right) =1$.}
\end{ex}

% nguồn: Ok_Toan11_TracNghiem_Lan2
\begin{ex}%[1D3H2-5]%[ma:47d467a9]
	Cho giới hạn $\lim\limits_{x\to +\infty } \left( \sqrt{36x^2+5ax+1}-6x+b \right)=\dfrac{20}{3}$ và đường thẳng $\Delta \colon y=ax+6b$ đi  qua điểm $M\left( 3;42 \right)$ với $a$, $b\in \mathbb{R}$. Giá trị của biểu thức $T=a^2+b^2$ là
	 \choice
		 {$104$}
		 {$100$}
		{\True $41$} 
		{$169$} 
	 \loigiai{
		Đường thẳng $\Delta\colon y=ax+6b$ đi  qua điểm $M\left( 3;42 \right)$ nên $3a+6b=42\Rightarrow a+2b=14$.\\
		Ta có $\lim\limits_{x\to +\infty  } \left( \sqrt{36x^2+5ax+1}-6x+b \right)=\lim\limits_{x\to +\infty } \left( \dfrac{5ax+1}{\sqrt{36x^2+5ax+1}+6}+b \right)$\\
		$=\lim\limits_{x\to  +\infty } \left( \dfrac{5a+\dfrac{1}{x}}{\sqrt{36+\dfrac{5a}{x}+\dfrac{1}{x^2}}+6}+b \right)=\dfrac{5a}{12}+b$.\\
		 Do đó $\dfrac{5a}{12}+b=\dfrac{20}{3}\Rightarrow 5a+12b=80$.\\
		Ta có hệ phương  trình $\heva{
		& 5a+12b=80\\ 
		 & a+2b=14\\
		}$  $\Leftrightarrow \heva{
		& a=4\\ 
		&  b=5.\\
		}$\\ 
	Vậy $T=a^2+b^2=41$.} 
\end{ex}

% nguồn: Ok_Toan11_TracNghiem_Lan2
\begin{ex}%[1D3H2-5]%[ma:cd88b076]
	Giới hạn $\lim\limits_{x\to 6} \dfrac{\sqrt{x-2}-2}{x-6}$  bằng
	 \choice
		{$\dfrac{1}{2}$} 
		 {$1$}
		 {\True $\dfrac{1}{4}$}
		 {$2$}
	\loigiai{ 
		Ta có $\lim\limits_{x\to 6} \dfrac{\sqrt{x-2}-2}{x-6}=\lim\limits_{x\to 6} \dfrac{x-6}{\left(x-6\right)\left(\sqrt{x-2}+2\right)}=\dfrac{1}{4}$. 
	 }
\end{ex}

% nguồn: Ok_Toan11_TracNghiem_Lan2
\begin{ex}%[1D3H2-5]%[ma:6ff7d070]
	 Tính $\lim\limits_{x\to +\infty } \left( \sqrt{x^2-4x+2}-x \right)$.
	 \choice
		{$-4$} 
		{\True $-2$} 
		{$4$} 
		 {$2$}
	\loigiai{ 
		Ta  có
		\begin{align*}
			&\lim\limits_{x\to +\infty  } \left( \sqrt{x^2-4x+2}-x \right) \\
			&= \lim\limits_{x\to +\infty } \dfrac{x^2-4x+2-x^2}{\sqrt{x^2-4x+2}+x} \\ 
			&= \lim\limits_{x\to +\infty } \dfrac{-4x+2}{\sqrt{x^2-4x+2}+x}  \\
			&= \lim\limits_{x\to +\infty } \dfrac{-4+\dfrac{2}{x}}{\sqrt{1-\dfrac{4}{x}+\dfrac{2}{x^2}}+1} \\ 
			&= -2. 
		\end{align*}
	} 
\end{ex}

% nguồn: Ok_Toan11_TracNghiem_Lan2
\begin{ex}%[1D3H2-5]%[ma:5d5fb7e9]
	Cho $\lim\limits_{x\to  -\infty } \left( \sqrt{x^2+ax+5}+x \right)=5$. Khi đó giá trị $a$ là
	\choice 
		 {$10$}
		 {$-6$}
		{$6$} 
		{\True $-10$} 
	 \loigiai{
		Ta có \begin{align*}
		&\lim\limits_{x\to -\infty  } \left( \sqrt{x^2+ax+5}+x \right) \\
		&= \lim\limits_{x\to  -\infty } \dfrac{\left( \sqrt{x^2+ax+5}+x \right)\left( \sqrt{x^2+ax+5}-x \right)}{\sqrt{x^2+ax+5}-x} \\
		&= \lim\limits_{x\to -\infty  } \dfrac{ax+5}{\sqrt{x^2+ax+5}-x} \\
		&= \lim\limits_{x\to -\infty  } \dfrac{a+\dfrac{5}{x}}{-\sqrt{1+\dfrac{a}{x}+\dfrac{5}{x^2}}-1} \\
		&= \dfrac{a}{-2}. 
	\end{align*}
Do đó $\lim\limits_{x\to -\infty } \left( \sqrt{x^2+ax+5}+x \right)=5\Leftrightarrow \dfrac{a}{-2}=5\Leftrightarrow a=-10$.} 
\end{ex}

% nguồn: Ok_Toan11_TracNghiem_Lan2
\begin{ex}%[1D3H2-5]%[ma:925734db]
	Biết $\lim\limits_{x\to-\infty}\left(\sqrt{5x^2+2x}+\sqrt{5} x\right)=\sqrt{5} a+b$  với $a,b\in\mathbb{Q}$. Tính $S=5a+b$.
	\choice 
		 {$S=-5$}
		{\True  $S=-1$}
		 {$S=1$}
		 {$S=5$}
	\loigiai{ 
	Ta  có
	$\lim\limits_{x\to-\infty}\left(\sqrt{5 x^2+2 x}+\sqrt{5}  x\right)=\lim\limits_{x\to-\infty} \dfrac{2 x}{\sqrt{5 x^2+2 x}-\sqrt{5} x}=\lim\limits_{x\to-\infty} \dfrac{2}{-\sqrt{5+\dfrac{2}{x}}-\sqrt{5}}=-\dfrac{\sqrt{5}}{5}$.\\
	Vậy  $\heva{&a=-\dfrac{1}{5}\\&b=0} \Rightarrow S=5 a+b=-1$.
	 }
\end{ex}

% nguồn: Ok_Toan11_TracNghiem_Lan2
\begin{ex}%[1D3H2-5]%[ma:d62f4f99]
	$\lim\limits_{x\to  +\infty } \left( \sqrt{x^2-5x+6}-x \right)$ bằng
	 \choice
		{$3$} 
		{$\dfrac{5}{2}$} 
		{\True  $-\dfrac{5}{2}$}
		{$-3$} 
	\loigiai{ 
		Ta có \begin{align*}
		&\lim\limits_{x\to  +\infty } \left( \sqrt{x^2-5x+6}-x \right)\\
		&=  \lim\limits_{x\to +\infty } \dfrac{-5x+6}{\sqrt{x^2-5x+6}+x} \\
		&= \lim\limits_{x\to +\infty } \dfrac{-5+\dfrac{6}{x}}{\sqrt{1-\dfrac{5}{x}+\dfrac{6}{x^2}}+1}  \\
		&=  -\dfrac{5}{2}.
	\end{align*}
} 
\end{ex}

% nguồn: Ok_Toan11_TracNghiem_Lan2
\begin{ex}%[1D3H2-5]%[ma:cc4f96fd]
	$\lim\limits_{x\to 0}\dfrac{\sqrt{1-x}-1}{x}$  bằng
	\choice 
		 {$\dfrac{1}{2}$}
		 {\True $-\dfrac{1}{2}$}
		{$0$} 
		 {$+\infty$}
	\loigiai{ 
		Ta có $\lim\limits_{x\to 0}\dfrac{\sqrt{1-x}-1}{x}=\lim\limits_{x\to 0}\dfrac{(\sqrt{1-x}-1)(\sqrt{1-x}+1)}{x(\sqrt{1-x}+1)}=\lim\limits_{x\to 0}\dfrac{1-x-1}{x(\sqrt{1-x}+1)}=-\dfrac{1}{2}$. 
	} 
\end{ex}

% nguồn: Ok_Toan11_TracNghiem_Lan2
\begin{ex}%[1D3H2-5]%[ma:40bb1105]
	Giá trị $\lim\limits_{x\to 3}  \dfrac{\sqrt{x+1}-2}{x-3}$ bằng
	 \choice
		{\True  $\dfrac{1}{4}$}
		{$\dfrac{1}{2}$} 
		 {$-\dfrac{1}{4}$}
		 {$-\dfrac{1}{2}$}
	 \loigiai{
		Ta có $\lim\limits_{x\to 3}  \dfrac{\sqrt{x+1}-2}{x-3}=\lim\limits_{x\to 3} \dfrac{x-3}{\left(\sqrt{x+1}+2\right)\left(x-3\right)}$ $=\lim\limits_{x\to 3} \dfrac{1}{\sqrt{x+1}+2}=\dfrac{1}{4}$.
	 }
\end{ex}

% nguồn: Ok_Toan11_TracNghiem_Lan2
\begin{ex}%[1D3H2-5]%[ma:a8e66b7b]
	Tính $\lim\limits_{x\to -1}\dfrac{\sqrt{5+x}-2}{x^2-1}$. 
	\choice 
		{$\dfrac{1}{8}$} 
		 {$-\dfrac{1}{4}$}
		{\True $-\dfrac{1}{8}$} 
		{$\dfrac{1}{4}$} 
	 \loigiai{Ta có
		\begin{align*}
			&\lim\limits_{x\to -1}\dfrac{\sqrt{5+x}-2}{x^2-1}\\ 
			=&\lim\limits_{x\to  -1}\dfrac{(\sqrt{5+x}-2)(\sqrt{5+x}+2)}{(x-1)(x+1)(\sqrt{5+x}+2)}\\
			 =&\lim\limits_{x\to -1}\dfrac{x+1}{(x-1)(x+1)(\sqrt{5+x}+2)}\\
			 =&-\dfrac{1}{8}.
		\end{align*}
	} 
\end{ex}

% nguồn: Ok_Toan11_TracNghiem_Lan2
\begin{ex}%[1D3H2-5]%[ma:01a0b031]
	Biết $\lim\limits_{x\to 0}\dfrac{\sqrt{3x+1}-1}{x}=\dfrac{a}{b}$, trong đó $a, b$ là các số nguyên dương và phân số $\dfrac{a}{b}$ tối giản. Tính giá trị biểu thức $P=a^2+b^2$. 
	 \choice
		{\True $P=13$} 
		 {$P=0$}
		{$P=5$} 
		 {$P=40$}
	\loigiai{ 
		Ta có $\lim\limits_{x\to 0}\dfrac{\sqrt{3x+1}-1}{x}=\lim\limits_{x\to 0}\dfrac{3x+1-1}{x(\sqrt{3x+1}+1)}=\lim\limits_{x\to  0}\dfrac{3}{\sqrt{3x+1}+1}=\dfrac{3}{2}$.\\
		Do đó $a=3$,  $b=2$.\\
		Vậy $P=a^2+b^2=13$. 
	} 
\end{ex}

% nguồn: Ok_Toan11_TracNghiem_Lan2
\begin{ex}%[1D3H2-5]%[ma:c413f415]
	Cho $\lim\limits_{x\to 2}\dfrac{\sqrt{2x+5}-3}{x-2}=\dfrac{a}{b}$, trong đó $a$, $b$ là các số nguyên dương và phân số $\dfrac{a}{b}$ tối giản. Tính giá  trị biểu thức $P=1984a^2+4b^2$.
	 \choice
		 {$0$}
		 {$2000$}
		{$8000$} 
		{\True $2020$} 
	 \loigiai{
		Ta  có
		\begin{align*}
			 &\lim\limits_{x\to 2}\dfrac{\sqrt{2x+5}-3}{x-2}\\
			=&\lim\limits_{x\to  2}\dfrac{(2x+5)-9}{(x-2)(\sqrt{2x+5}+3)}\\
			 =&\lim\limits_{x\to 2}\dfrac{2(x-2)}{(x-2)(\sqrt{2x+5}+3)}\\
			=&\lim\limits_{x\to 2}\dfrac{2}{\sqrt{2x+5}+3}\\ 
			 =&\dfrac{1}{3}\\
			\Rightarrow&\heva{&a=1\\&b=3.} 
		\end{align*}
		 Vậy $P=1984\cdot1^2+4\cdot3^2=2020$.
	} 
\end{ex}

% nguồn: Ok_Toan11_TracNghiem_Lan2
\begin{ex}%[1D3H2-5]%[ma:b3effa81]
	Tính giới hạn $\lim\limits_{x\to 0} \dfrac{\sqrt{4x^2-2x+1}-\sqrt{1-2x}}{x}$. 
	 \choice
		 {$2$}
		{$-1$} 
		 {$-2$}
		{\True  $0$}
	 \loigiai{
		Ta  có
		\begin{eqnarray*}
			\lim\limits_{x\to 0} \dfrac{\sqrt{4x^2-2x+1}-\sqrt{1-2x}}{x}&=&\lim\limits_{x\to  0} \dfrac{4x^2}{x\left(\sqrt{4x^2-2x+1}+\sqrt{1-2x} \right)}\\
			&=& \lim\limits_{x\to 0} \dfrac{4x}{\left(\sqrt{4x^2-2x+1}+\sqrt{1-2x}  \right)}=0.
		\end{eqnarray*}
	 }
\end{ex}

% nguồn: Ok_Toan11_TracNghiem_Lan2
\begin{ex}%[1D3H2-5]%[ma:9897ab3d]
	Số nào trong  các số sau là bằng $\lim\limits_{x\to 3}\dfrac{\sqrt{x^2+x}-2\sqrt{3}}{x-3}$?
	 \choice
		 {$\dfrac{\sqrt{3}}{12}$}
		{$-\dfrac{\sqrt{3}}{12}$} 
		 {\True $\dfrac{7\sqrt{3}}{12}$}
		 {$-\dfrac{7\sqrt{3}}{12}$}
	\loigiai{ 
		Ta  có
		\begin{align*}
			&\lim\limits_{x\to 3}\dfrac{\sqrt{x^2+x}-2\sqrt{3}}{x-3}\\ 
			=&\lim\limits_{x\to  3}\dfrac{x^2+x-12}{(x-3)(\sqrt{x^2+x}+2\sqrt{3})}\\ =&\lim\limits_{x\to 3}\dfrac{(x-3)(x+4)}{(x-3)(\sqrt{x^2+x}+2\sqrt{3})}\\
			=&\lim\limits_{x\to 3}\dfrac{x+4}{\sqrt{x^2+x}+2\sqrt{3}}\\ 
			=&\dfrac{3+4}{\sqrt{3^2+3}+2\sqrt{3}}\\ 
			=&\dfrac{7}{4\sqrt{3}}=\dfrac{7\sqrt{3}}{12}. 
		\end{align*}
	 }
\end{ex}

% nguồn: Ok_Toan11_TracNghiem_Lan2
\begin{ex}%[1D3H2-5]%[ma:bae61433]
	$\lim\limits_{x\to  -2} \dfrac{x^2-2x-8}{\sqrt{2x+5}-1}$ bằng
	 \choice
		{$-3$} 
		 {$\dfrac{1}{\sqrt{2}}$}
		 {\True $-6$}
		{$8$} 
	\loigiai{ 
		 Ta có
		\begin{align*}
			\lim\limits_{x\to -2} \dfrac{x^2-2x-8}{\sqrt{2x+5}-1} &=\lim\limits_{x\to  -2} \dfrac{\left(x^2-2x-8\right)\left(\sqrt{2x+5}+1\right)}{\left(\sqrt{2x+5}-1\right)\left(\sqrt{2x+5}+1\right)}\\
			&=\lim\limits_{x\to -2}  \dfrac{\left(x+2\right)\left(x-4\right)\left(\sqrt{2x+5}+1\right)}{2\left(x+2\right)}\\
			 &=\lim\limits_{x\to -2} \dfrac{\left(x-4\right)\left(\sqrt{2x+5}+1\right)}{2}\\
			&=\dfrac{-6\cdot  (1+1)}{2}=-6.
		\end{align*}
		Vậy  $\lim\limits_{x\to -2} \dfrac{x^2-2x-8}{\sqrt{2x+5}-1}=-6$.
	} 
\end{ex}

% nguồn: Ok_Toan11_TracNghiem_Lan2
\begin{ex}%[1D3H2-5]%[ma:b06208e7]
	Tính 
	$I=\lim\limits_{x\to+\infty}\left(\sqrt{x^2-4  x+2}-x\right)$.
	 \choice
		 {$I=-4$}
		{\True  $I=-2$}
		{$I=4$} 
		{$I=2$} 
	 \loigiai{
	Ta có $\begin{aligned}
	I&=\lim\limits_{x\to+\infty}\left(\sqrt{x^2-4 x+2}-x\right)=\lim\limits_{x\to+\infty} \dfrac{x^2-4  x+2-x^2}{\sqrt{x^2-4 x+2}+x}\\&=\lim\limits_{x\to+\infty} \dfrac{-4 x+2}{\sqrt{x^2-4 x+2}+x}=\lim\limits_{x\to+\infty} \dfrac{-4+\dfrac{2}{x}}{\sqrt{1-\dfrac{4}{x}+\dfrac{2}{x^2}}+1}=-2.
\end{aligned}$
} 
\end{ex}

% nguồn: Ok_Toan11_TracNghiem_Lan2
\begin{ex}%[1D3H2-5]%[ma:587f66e3]
	Giới hạn $\lim\limits_{x\to 0}\dfrac{\sqrt{5x+3}-\sqrt{3}}{x}=\dfrac{m}{n\sqrt{k}}$  ($m$, $n$, $k$ nguyên dương và nguyên tố cùng nhau). Tính $m-n+k$.
	\choice 
		 {\True $6$}
		{$4$} 
		{$8$} 
		{$0$} 
	 \loigiai{
		Ta có $\lim\limits_{x\to 0}\dfrac{\sqrt{5x+3}-\sqrt{3}}{x}=\lim\limits_{x\to 0}\dfrac{5x}{x(\sqrt{5x+3}+\sqrt{3})}=\dfrac{5}{2\sqrt{3}}$.\\ 
		Suy ra  $m=5$; $n=2$; $k=3$ nên $m-n+k=6$.
	} 
\end{ex}

% nguồn: Ok_Toan11_TracNghiem_Lan2
\begin{ex}%[1D3H2-5]%[ma:fce8bf9f]
	Giới hạn $\lim\limits_{x\to  5}\dfrac{\sqrt{x-1}-2}{x-5}$ bằng
	\choice 
		{\True  $\dfrac{1}{4}$}
		 {$2$}
		{$1$} 
		 {$\dfrac{1}{2}$}
	 \loigiai{ Ta có
		 $\lim\limits_{x\to 5}\dfrac{\sqrt{x-1}-2}{x-5}=\lim\limits_{x\to 5}\dfrac{x-5}{(x-5)(\sqrt{x-1}+2)}=\lim\limits_{x\to 5}\dfrac{1}{\sqrt{x-1}+2}=\dfrac{1}{4}$.
	 }
\end{ex}

% nguồn: Ok_Toan11_TracNghiem_Lan2
\begin{ex}%[1D3H2-5]%[ma:6e263f22]
	Cặp $(a;b)$ thỏa mãn $\lim\limits_{x\to 3}\dfrac{x^2+ax+b}{x-3}=3$ là 
	\choice 
		{\True $a=-3$, $b=0$} 
		{$a=3$,  $b=0$}
		 {$a=0$, $b=-9$}
		{Không tồn tại cặp $(a;b)$ thỏa mãn} 
	\loigiai{ 
		Để $\lim\limits_{x\to 3}\dfrac{x^2+ax+b}{x-3}=3$ thì  ta phải có $x^2+ax+b=(x-3)(x-m)$.\\
		Khi  đó $3-m=3\Leftrightarrow m=0$.\\
		Vậy  $x^2+ax+b=(x-3)x=x^2-3x$.\\
		Suy ra $a=-3$ và $b=0$. 
	} 
\end{ex}

% nguồn: Ok_Toan11_TracNghiem_Lan2
\begin{ex}%[1D3H2-5]%[ma:c486b261]
	 Trong các giới hạn sau, giới hạn nào có kết quả là $0$?
	\choice 
		{$\lim\limits_{x\to 1} \dfrac{x-1}{x^3-1}$} 
		{$\lim\limits_{x\to  -2} \dfrac{2x+5}{x+10}$}
		{$\lim\limits_{x\to  1} \dfrac{x^2-1}{x^2-3x+2}$}
		{\True $\lim\limits_{x\to +\infty }  \left( \sqrt{x^2+1}-x \right)$}
	 \loigiai{
	Xét $\lim\limits_{x\to +\infty } \left(  \sqrt{x^2+1}-x \right)=\lim\limits_{x\to +\infty } \dfrac{x^2+1-x^2}{\sqrt{x^2+1}+x}=\lim\limits_{x\to +\infty } \dfrac{1}{\sqrt{x^2+1}+x}=0$.}
\end{ex}

% nguồn: Ok_Toan11_TracNghiem_Lan2
\begin{ex}%[1D3H2-5]%[ma:a8de9757]
	Tính  giới hạn $K=\lim\limits_{x\to 0} \dfrac{\sqrt{4x+1}-1}{x^2-3x}$.
	\choice 
		{\True  $K=-\dfrac{2}{3}$}
		{$K=\dfrac{2}{3}$} 
		{$K=\dfrac{4}{3}$} 
		{$K=0$} 
	 \loigiai{
		Ta có $K=\lim\limits_{x\to 0} \dfrac{\sqrt{4x+1}-1}{x^2-3x}=\lim\limits_{x\to 0} \dfrac{4x}{x\left(x-3\right)\left(\sqrt{4x+1}+1\right)}=\lim\limits_{x\to 0}  \dfrac{4}{\left(x-3\right)\left(\sqrt{4x+1}+1\right)}=-\dfrac{2}{3}$.
	} 
\end{ex}

% nguồn: Ok_Toan11_TracNghiem_Lan2
\begin{ex}%[1D3H2-5]%[ma:0183e958]
	Tính $\lim\limits_{x\to 1}\dfrac{\sqrt{x^2+x+2}-2}{x-1}$. 
	 \choice
		{$-\dfrac{1}{4}$} 
		{\True  $\dfrac{3}{4}$}
		 {$1$}
		{$-\infty$} 
	 \loigiai{
		Ta có $\lim\limits_{x\to 1}\dfrac{\sqrt{x^2+x+2}-2}{x-1}=\lim\limits_{x\to 1}\dfrac{x^2+x-2}{(x-1)(\sqrt{x^2+x+2}+2)}=\lim\limits_{x\to 1}\dfrac{(x-1)(x+2)}{(x-1)(\sqrt{x^2+x+2}+2)}=\dfrac{3}{4}$. 
	 }
\end{ex}

% nguồn: Ok_Toan11_TracNghiem_Lan2
\begin{ex}%[1D3H2-5]%[ma:77df2d53]
	Trong bốn giới hạn sau đây,  giới hạn nào bằng $0$?
	\choice 
		 {$\lim\limits_{x\to 1} \dfrac{x^2-1}{3x^2-4x+1}$}
		{$\lim\limits_{x\to -2} \dfrac{4x+3}{x+5}$} 
		 {$\lim\limits_{x\to 1} \dfrac{x^2-1}{x^3-1}$}
		{\True $\lim\limits_{x\to+\infty} \left(\sqrt{x^2+1}-x\right)$} 
	 \loigiai{
		\begin{itemize}
			 \item $\lim\limits_{x\to 1} \dfrac{x^2-1}{3x^2-4x+1}=\lim\limits_{x\to 1} \dfrac{x+1}{3x-1}=\dfrac{2}{2}=1$.
			 \item $\lim\limits_{x\to -2} \dfrac{4x+3}{x+5}=\dfrac{-5}{3}$.
			\item $\lim\limits_{x\to 1} \dfrac{x^2-1}{x^3-1}=\lim\limits_{x\to 1} \dfrac{x+1}{x^2+x+1}=\dfrac{2}{3}$. 
			\item $\lim\limits_{x\to+\infty} \left(\sqrt{x^2+1}-x\right)=\lim\limits_{x\to+\infty} \dfrac{1}{\sqrt{x^2+1}+x}=0$. 
		\end{itemize}
	} 
\end{ex}

% nguồn: Ok_TracNghiemLop11_Lan1
\begin{ex}%[1D3H2-5]%[ma:39d45db5]
	Giá trị của giới  hạn $\lim\limits_{x\to+\infty}\left(\sqrt{x^2+1}-x\right)$ là
	\choice 
		{\True  $0$}
		{$+\infty $} 
		 {$\dfrac{1}{2}$}
		{$-\infty  $}
	 \loigiai{
		$\lim\limits_{x\to+\infty}\left(\sqrt{x^2+1}-x\right)= \lim\limits_{x\to+\infty}\dfrac{1}{\sqrt{x^2+1}+x}=  \lim\limits_{x\to+\infty}\dfrac{\dfrac{1}{x}}{\sqrt{1+\dfrac{1}{x^2}}+1}= \dfrac{0}{2}=0$.
	} 
\end{ex}

% nguồn: Ok_TracNghiemLop11_Lan1
\begin{ex}%[1D3H2-5]%[ma:e195c16f]
	Tính  $\displaystyle\lim\limits_{x\to 0} \dfrac{\sqrt{2x+1}-1}{x}$.
	 \choice
		{\True  $1$}
		 {$5$}
		 {$2$}
		{$4$} 
	 \loigiai{
		Ta có $\displaystyle\lim\limits_{x\to 0} \dfrac{\sqrt{2x+1}-1}{x}=\displaystyle\lim\limits_{x\to  0}\dfrac{2x}{x\left(\sqrt{2x+1}+1\right)}=\displaystyle\lim\limits_{x\to 0}\dfrac{2}{\sqrt{2x+1}+1}=1$.
	 }
\end{ex}

% nguồn: Ok_TracNghiemLop11_Lan1
\begin{ex}%[1D3H2-5]%[ma:c197a758]
	Tìm giới hạn sau $\lim\limits_{x \to 5} \dfrac{\sqrt{x + 76}  - 9}{x - 5}$.
	 \choice
		 {\True $\dfrac{1}{18}$}
		{$-  \dfrac{17}{18}$}
		 {$\dfrac{19}{18}$}
		{$\dfrac{5}{9}$} 
	 \loigiai{
		Nhân  cả tử và mẫu với $\sqrt{x + 76} + 9$, ta được
		\begin{eqnarray*}
			\lim\limits_{x \to 5} \dfrac{(\sqrt{x + 76} - 9)(\sqrt{x + 76} + 9)}{(x - 5)(\sqrt{x + 76} + 9)} &=& \lim\limits_{x  \to 5} \dfrac{x + 76 - 9^2}{(x - 5)(\sqrt{x + 76} + 9)}. \\
			&=& \lim\limits_{x \to 5} \dfrac{(x  - 5)}{(x - 5)(\sqrt{x + 76} + 9)}\\
			&=& \dfrac{1}{\sqrt{81} + 9} =  \dfrac{1}{18}.
		\end{eqnarray*}
	Vậy $\lim\limits_{x \to 5} \dfrac{\sqrt{x  + 76} - 9}{x - 5} = \dfrac{1}{18}.$}
\end{ex}

% nguồn: Ok_TracNghiemLop11_Lan1
\begin{ex}%[1D3H2-5]%[ma:1b40b256]
	Giá trị của  giới hạn $\lim\limits_{x\to+\infty}(\sqrt{x^2+3x}-\sqrt{x^2+4x})$ là
	 \choice
		 {$\dfrac{7}{2}$}
		{\True  $-\dfrac{1}{2}$}
		{$+\infty  $}
		{$-\infty  $}
	\loigiai{ 
		\begin{eqnarray*}
			\lim\limits_{x\to+\infty}(\sqrt{x^2+3x}-\sqrt{x^2+4x}) &=&\lim\limits_{x\to+\infty}\dfrac{-x}{\sqrt{x^2+3x}+\sqrt{x^2+4x}}\\ 
			&=&\lim\limits_{x\to+\infty}\dfrac{-1}{\sqrt{1+\dfrac{3}{x}}+\sqrt{1+\dfrac{4}{x}}}=-\dfrac{1}{2}. 
		\end{eqnarray*}
	} 
\end{ex}

% nguồn: Ok_Toan11_TracNghiem_Lan2
\begin{ex}%[1D3V2-5]%[ma:5627abb4]
	Kết  quả của $\lim\limits_{x\to-\infty}\dfrac{1}{x}\left(\sqrt{x^2+16}\right)$ bằng
	 \choice
		 {\True $-1$}
		 {$-4$}
		{$-\infty$} 
		{$+\infty$} 
	\loigiai{ 
		 Ta có
		 \allowdisplaybreaks
		\begin{align*}
			 \lim\limits_{x\to-\infty}\dfrac{1}{x}\sqrt{x^2+16} &=\lim\limits_{x\to-\infty}\dfrac{1}{x}\cdot |x|\sqrt{1+\dfrac{16}{x^2}}=\lim\limits_{x\to-\infty}\left(-\sqrt{1+\dfrac{16}{x^2}}\right)=-1.
		\end{align*}
		Vậy  $\lim\limits_{x\to-\infty}\dfrac{1}{x}\left(\sqrt{x^2+16}\right)=-1.$
	} 
\end{ex}

% nguồn: Ok_Toan11_TracNghiem_Lan2
\begin{ex}%[1D3V2-5]%[ma:277491d7]
	Cho $\lim\limits_{x\to+\infty}\left(\sqrt[3]{8x^3-ax^2+10}-2bx\right)=1$, khi đó 
	\choice 
		 {$2a+b=-11$}
		{\True $a+b=-11$} 
		{$a+2b=-11$} 
		{$a-b=-11$} 
	 \loigiai{
		Ta có  $b=1$ vì nếu $b\ne 1$ thì
		\begin{align*}
			\lim\limits_{x\to+\infty}\left(\sqrt[3]{8x^3-ax^2+10}-2bx\right) 
			 =\lim\limits_{x\to+\infty} x\left(\sqrt[3]{8-\dfrac{a}{x}+\dfrac{1}{x^3}}-2b\right)=\pm\infty.
		\end{align*}
		 Khi đó
		\allowdisplaybreaks 
		\begin{align*}
			 \lim\limits_{x\to+\infty}\left(\sqrt[3]{8x^3-ax^2+10}-2x\right)
			&=\lim\limits_{x\to+\infty}\dfrac{8x^3-ax^2+10-8x^3}{\sqrt[3]{(8x^3-ax^2+10)^2}+2x\sqrt[3]{8x^3-ax^2+10}+4x^2}\\ 
			&=\lim\limits_{x\to+\infty}\dfrac{-ax^2+10}{\sqrt[3]{(8x^3-ax^2+10)^2}+2x\sqrt[3]{8x^3-ax^2+10}+4x^2}\\ 
			&=\lim\limits_{x\to+\infty}\dfrac{-a+\dfrac{10}{x^2}}{\sqrt[3]{\left(8-\dfrac{a}{x}+\dfrac{1}{x^3}\right)^2}+2\sqrt[3]{8-\dfrac{a}{x}+\dfrac{1}{x^3}}+4}\\ 
			 &=\dfrac{-a}{12}=1.
		\end{align*}
		Nên $a=-12$. Vậy  $a+b=-11$.
	 }
\end{ex}

% nguồn: Ok_Toan11_TracNghiem_Lan2
\begin{ex}%[1D3V2-5]%[ma:0368b824]
	 Kết quả của $\lim\limits_{x\to+\infty} (1+x)\sqrt{\dfrac{3x}{x^2-1}}$ bằng
	\choice 
		 {$-2\sqrt{3}$}
		 {$2\sqrt{3}$}
		{$-\infty$} 
		 {\True $+\infty$}
	 \loigiai{
		Khi $x\to+\infty$ thì $1+x > 0$ nên ta  có
		\allowdisplaybreaks 
		\begin{align*}
			 \lim\limits_{x\to+\infty} (1+x)\sqrt{\dfrac{3x}{x^2-1}} &=\lim\limits_{x\to+\infty}\sqrt{\dfrac{3x(1+x)^2}{x^2-1}}\\
			 &=\lim\limits_{x\to+\infty}\sqrt{\dfrac{\left(\dfrac{1}{x}+1\right)^2\cdot 3x}{1-\dfrac{1}{x^2}}}=+\infty.
		\end{align*}
		Vậy $\lim\limits_{x\to+\infty}  (1+x)\sqrt{\dfrac{3x}{x^2-1}}=+\infty$.
	 }
\end{ex}

% nguồn: Ok_Toan11_TracNghiem_Lan2
\begin{ex}%[1D3V2-5]%[ma:98769850]
	Tính  giới hạn $\lim\limits_{x\to+\infty} \left(\sqrt{x^2-4x+2}-x\right)$.
	 \choice
		 {$-4$}
		 {\True $-2$}
		 {$4$}
		{$2$} 
	\loigiai{ 
		Ta có \begin{eqnarray*}
		\lim\limits_{x\to+\infty}  \left(\sqrt{x^2-4x+2}-x\right)&=&\lim\limits_{x\to+\infty} \dfrac{x^2-4x+2-x^2}{\sqrt{x^2-4x+2}+x}\\
		 &=&\lim\limits_{x\to+\infty} \dfrac{-4x+2}{\sqrt{x^2-4x+2}+x}\\
		 &=&\lim\limits_{x\to+\infty} \dfrac{x\left(-4+\dfrac{2}{x} \right)}{x\left(\sqrt{1-\dfrac{4}{x}+\dfrac{2}{x^2}}+1\right)}\\
		 &=&\lim\limits_{x\to+\infty} \dfrac{-4+\dfrac{2}{x}}{\sqrt{1-\dfrac{4}{x}+\dfrac{2}{x^2}}+1}\\
		&=&-2. 
	\end{eqnarray*}
 }
\end{ex}

% nguồn: Ok_Toan11_TracNghiem_Lan2
\begin{ex}%[1D3V2-5]%[ma:0abcab76]
	 Kết quả của $\lim\limits_{x\to+\infty}\dfrac{1}{x}\left(\sqrt{x^2+16}-x\right)$ bằng
	 \choice
		{$-\infty$} 
		{$4$} 
		{\True $0$} 
		 {$+\infty$}
	\loigiai{ 
		 Ta có
		\allowdisplaybreaks 
		\begin{align*}
			\lim\limits_{x\to+\infty}\dfrac{1}{x}\left(\sqrt{x^2+16}-x\right) &=\lim\limits_{x\to+\infty}\dfrac{1}{x}\left(\dfrac{x^2+16-x^2}{\sqrt{x^2+16}+x}\right)\\ 
			 &=\lim\limits_{x\to+\infty}\dfrac{1}{x}\left(\dfrac{16}{\sqrt{x^2+16}+x}\right)=0.
		\end{align*}
	 Vậy $\lim\limits_{x\to+\infty}\dfrac{1}{x}\left(\sqrt{x^2+16}-x\right)=0$.}
\end{ex}

% nguồn: Ok_Toan11_TracNghiem_Lan2
\begin{ex}%[1D3V2-5]%[ma:3ca7499b]
	 Giới hạn $\lim\limits_{x\to+\infty} x\left(\sqrt{x^2-1}-x\right)$ bằng
	 \choice
		{\True $-\dfrac{1}{2}$} 
		{$1$} 
		{$+\infty$} 
		 {$0$}
	 \loigiai{
		\begin{eqnarray*}
			\lim\limits_{x\to+\infty} x\left(\sqrt{x^2-1}-x\right)&=&\lim\limits_{x\to+\infty} \dfrac{-x}{\sqrt{x^2-1}+x}\\ 
			&=&\lim\limits_{x\to+\infty}  \dfrac{-1}{\sqrt{1-\dfrac{1}{x^2}}+1}\\
			 &=&-\dfrac{1}{2}.
		\end{eqnarray*}
		Vậy  $\lim\limits_{x\to+\infty} x\left(\sqrt{x^2-1}-x\right)=-\dfrac{1}{2}$.
	 }
\end{ex}

% nguồn: Ok_Toan11_TracNghiem_Lan2
\begin{ex}%[1D3V2-5]%[ma:c03448a9]
	Giới hạn $\lim\limits_{x\to+\infty} \left(x-\sqrt{x^2-4x} \right)$ bằng 
	\choice 
		{$-2$} 
		 {$4$}
		 {$-4$}
		 {\True $2$}
	 \loigiai{
		Ta  có
		\begin{eqnarray*}
			\lim\limits_{x\to+\infty} \left(x-\sqrt{x^2-4x}  \right)&=&\lim\limits_{x\to+\infty} \dfrac{x^2-\left(x^2-4x\right)}{x+\sqrt{x^2-4x}}\\
			 &=&\lim\limits_{x\to+\infty} \dfrac{4x}{x+\sqrt{x^2-4x}}\\
			&=&\lim\limits_{x\to+\infty}  \dfrac{4x}{x\left(1+\sqrt{1-\dfrac{4}{x}} \right)}\\
			&=&\lim\limits_{x\to+\infty} \dfrac{4}{1+\sqrt{1-\dfrac{4}{x}}}\\ 
			&=&2. 
		\end{eqnarray*}
		 Vậy $\lim\limits_{x\to+\infty} \left(x-\sqrt{x^2-4x} \right)=2$.
	} 
\end{ex}

% nguồn: Ok_Toan11_TracNghiem_Lan2
\begin{ex}%[1D3V2-5]%[ma:6c1af9e8]
	Cho $\lim\limits_{x\to-\infty} \left(\sqrt{x^2+ax+5}+bx\right)=5$. Khi đó giá trị  của $a+b$ là
	 \choice
		{\True $-9$} 
		 {$6$}
		{$9$} 
		 {$-6$}
	\loigiai{ 
		\begin{itemize}
			\item Nếu $b\le 0$ thì \begin{eqnarray*}
			\lim\limits_{x\to-\infty} \left(\sqrt{x^2+ax+5}+bx\right)&=&\lim\limits_{x\to-\infty}  x\left(-\sqrt{1+\dfrac{a}{x}+\dfrac{5}{x^2}}+b\right)\\
			 &=&+\infty.
		\end{eqnarray*}
		\item  Nếu $b=1$ thì
		\begin{eqnarray*}
			\lim\limits_{x\to-\infty}  \left(\sqrt{x^2+ax+5}+x\right)&=&\lim\limits_{x\to-\infty} \dfrac{ax+5}{\sqrt{x^2+ax+5}-x}\\
			&=&\lim\limits_{x\to-\infty} \dfrac{a+\dfrac{5}{x}}{-\sqrt{1+\dfrac{a}{x}+\dfrac{5}{x^2}}-1}\\ 
			 &=&\dfrac{a}{-2}\\
			 &=&5.
		\end{eqnarray*}
		Suy ra $a=-10$. Khi đó $a+b=-9$. 
		\item Nếu $\heva{&b\ne 1\\&b > 0}$ thì \begin{eqnarray*}
		 \lim\limits_{x\to-\infty} \left(\sqrt{x^2+ax+5}+bx\right)&=&\lim\limits_{x\to-\infty} \dfrac{\left(1-b^2\right)x^2+ax+5}{\sqrt{x^2+ax+5}-bx}\\
		&=&\pm \infty. 
	\end{eqnarray*}
	\end{itemize}
} 
\end{ex}

% nguồn: Ok_Toan11_TracNghiem_Lan2
\begin{ex}%[1D3V2-5]%[ma:cf1a8159]
	Kết quả của  $\lim\limits_{x\to 3^+} (x-3)\left(\dfrac{1}{\sqrt{x^2-9}}\right)$ bằng
	 \choice
		{$1$} 
		 {\True $0$}
		{$-\infty$} 
		 {$+\infty$}
	\loigiai{ 
		Ta có 
		\begin{align*}
			\lim\limits_{x\to  3^+} (x-3)\left(\dfrac{1}{\sqrt{x^2-9}}\right) &=\lim\limits_{x\to 3^+}\sqrt{x-3}\left(\dfrac{1}{\sqrt{x+3}}\right)=0.
		\end{align*}
		 Vậy $\lim\limits_{x\to 3^+} (x-3)\left(\dfrac{1}{\sqrt{x^2-9}}\right)=0$.
	 }
\end{ex}

% nguồn: Ok_Toan11_TracNghiem_Lan2
\begin{ex}%[1D3V2-5]%[ma:259b4951]
	Giới hạn nào dưới đây có kết quả  là $\dfrac{1}{2}$?
	 \choice
		 {$\lim\limits_{x\to-\infty} \dfrac{x}{2} \left(\sqrt{x^2+1}-x\right)$}
		 {$\lim\limits_{x\to+\infty} x\left(\sqrt{x^2+1}+x\right)$}
		{$\lim\limits_{x\to-\infty} \dfrac{x}{2} \left(\sqrt{x^2+1}+x\right)$} 
		{\True  $\lim\limits_{x\to+\infty} x\left(\sqrt{x^2+1}-x\right)$}
	\loigiai{ 
		Xét 
		\begin{eqnarray*}
			\lim\limits_{x\to+\infty}  x\left(\sqrt{x^2+1}-x\right)&=&\lim\limits_{x\to+\infty} \dfrac{x}{\sqrt{x^2+1}+x}\\
			&=&\lim\limits_{x\to+\infty}  \dfrac{x}{\left|x\right|\sqrt{1+\dfrac{1}{x^2}}+x}\\
			&=&\lim\limits_{x\to+\infty}  \dfrac{x}{x\sqrt{1+\dfrac{1}{x^2}}+x}\\
			&=&\lim\limits_{x\to+\infty}  \dfrac{1}{\sqrt{1+\dfrac{1}{x^2}}+1}\\
			 &=&\dfrac{1}{2}.
		\end{eqnarray*}
	} 
\end{ex}

% nguồn: Ok_Toan11_TracNghiem_Lan2
\begin{ex}%[1D3V2-5]%[ma:97d18f3b]
	Tìm  giới hạn $I=\lim\limits_{x\to+\infty} \left(x+1-\sqrt{x^2-x+2} \right)$.
	 \choice
		{$I=\dfrac{1}{2}$} 
		 {$I=\dfrac{46}{31}$}
		 {$I=\dfrac{17}{11}$}
		{\True  $I=\dfrac{3}{2}$}
	 \loigiai{
		Ta  có
		\begin{eqnarray*}
			I&=&\lim\limits_{x\to+\infty}  \left(x+1-\sqrt{x^2-x+2} \right)\\
			&=&\lim\limits_{x\to+\infty} \dfrac{x^2-x^2+x-2}{x+\sqrt{x^2-x+2}}+1\\ 
			 &=&\lim\limits_{x\to+\infty} \dfrac{x-2}{x+\sqrt{x^2-x+2}}+1\\
			 &=&\lim\limits_{x\to+\infty} \dfrac{1-\dfrac{2}{x}}{1+\sqrt{1-\dfrac{1}{x}+\dfrac{2}{x^2}}}+1\\
			&=&\dfrac{3}{2}. 
		\end{eqnarray*}
	} 
\end{ex}

% nguồn: Ok_Toan11_TracNghiem_Lan2
\begin{ex}%[1D3V2-5]%[ma:59a7ccaa]
	 Kết quả của $\lim\limits_{x\to+\infty} x\sqrt{\dfrac{2x+1}{3x^3+x^2+2}}$ bằng
	\choice 
		 {$\dfrac{2}{3}$}
		 {\True $\dfrac{\sqrt{6}}{3}$}
		{$-\infty$} 
		 {$+\infty$}
	\loigiai{ 
		 Ta có
		 \allowdisplaybreaks
		\begin{align*}
			\lim\limits_{x\to+\infty}  x\sqrt{\dfrac{2x+1}{3x^3+x^2+2}} &=\lim\limits_{x\to+\infty}\sqrt{\dfrac{(2x+1)x^2}{3x^3+x^2+2}}\\
			&=\lim\limits_{x\to+\infty}\sqrt{\dfrac{2x^3+x^2}{3x^3+x^2+2}}\\ 
			&=\lim\limits_{x\to+\infty}\sqrt{\dfrac{2+\dfrac{1}{x}}{3+\dfrac{1}{x}+\dfrac{2}{x^3}}}=\dfrac{\sqrt{6}}{3}. 
		\end{align*}
		Vậy  $\lim\limits_{x\to+\infty} x\sqrt{\dfrac{2x+1}{3x^3+x^2+2}}=\dfrac{\sqrt{6}}{3} $.
	 }
\end{ex}

% nguồn: Ok_Toan11_TracNghiem_Lan2
\begin{ex}%[1D3V2-5]%[ma:b4809a0f]
	Giới hạn $\lim\limits_{x\to+\infty} \left(\sqrt{x+\sqrt{3x+\sqrt{2x}}}-\sqrt{x} \right)$  bằng
	\choice 
		{$\dfrac{1}{2\sqrt{3}}$} 
		{$+\infty$} 
		{\True  $\dfrac{\sqrt{3}}{2}$}
		{$\dfrac{1}{\sqrt{2}}$} 
	 \loigiai{
		\begin{eqnarray*}
			\lim\limits_{x\to+\infty}  \left(\sqrt{x+\sqrt{3x+\sqrt{2x}}}-\sqrt{x} \right)&=&\lim\limits_{x\to+\infty} \dfrac{x+\sqrt{3x+\sqrt{2x}}-x}{\sqrt{x+\sqrt{3x+\sqrt{2x}}}+\sqrt{x}}\\
			 &=&\lim\limits_{x\to+\infty} \dfrac{\sqrt{3x+\sqrt{2x}}}{\sqrt{x+\sqrt{3x+\sqrt{2x}}}+\sqrt{x}}\\
			 &=&\lim\limits_{x\to+\infty} \dfrac{\sqrt{3+\sqrt{\dfrac{2}{x}}}}{\sqrt{1+\sqrt{\dfrac{3}{x}+\sqrt{\dfrac{2}{x^3}}}}+1} \\
			&=&\dfrac{\sqrt{3}}{2}. 
		\end{eqnarray*}
		Vậy  $\lim\limits_{x\to+\infty} \left(\sqrt{x+\sqrt{3x+\sqrt{2x}}}-\sqrt{x} \right)=\dfrac{\sqrt{3}}{2}$.
	 }
\end{ex}

% nguồn: Ok_Toan11_TracNghiem_Lan2
\begin{ex}%[1D3V2-5]%[ma:5b3873fe]
	Giá trị của giới hạn $\lim\limits_{x\to+\infty}  \left(\sqrt{x^2+3x}-\sqrt{x^2+4x} \right)$ là
	 \choice
		{$\dfrac{7}{2}$} 
		{\True  $-\dfrac{1}{2}$}
		 {$+\infty$}
		{$-\infty$} 
	\loigiai{ 
		Ta có \begin{eqnarray*}
		\lim\limits_{x\to+\infty}  \left(\sqrt{x^2+3x}-\sqrt{x^2+4x} \right)&=&\lim\limits_{x\to+\infty} \dfrac{-x}{\sqrt{x^2+3x}+\sqrt{x^2+4x}}\\
		 &=&\lim\limits_{x\to+\infty} \dfrac{-1}{\sqrt{1+\dfrac{3}{x}}+\sqrt{1+\dfrac{4}{x}}}\\
		&=&-\dfrac{1}{2}. 
	\end{eqnarray*}
} 
\end{ex}

% nguồn: Ok_Toan11_TracNghiem_Lan2
\begin{ex}%[1D3V2-5]%[ma:9f1698d1]
	Gọi $a,b$ là  các số thực thỏa mãn $\lim\limits_{x\to+\infty} \left(\sqrt{4x^2-3x+1}-\left(ax+b\right)\right)=0$. Khi đó $3a+8b$ bằng
	 \choice
		{$3$} 
		 {$5$}
		 {\True $0$}
		{$-1$} 
	\loigiai{ 
		Ta  có
		\begin{eqnarray*}
			\lim\limits_{x\to+\infty} \left(\sqrt{4x^2-3x+1}-\left(ax+b\right)\right)&=&\lim\limits_{x\to+\infty}  \dfrac{4x^2-3x+1-\left(ax+b\right)^2}{\sqrt{4x^2-3x+1}+ax+b}\\
			 &=&\lim\limits_{x\to+\infty} \dfrac{\left(4-a^2 \right)x^2-\left(3+2ab\right)x+1-b^2}{\sqrt{4x^2-3x+1}+ax+b}\\
			&=&\lim\limits_{x\to+\infty} \dfrac{\left(4-a^2  \right)x-\left(3+2ab\right)+\dfrac{1-b^2}{x}}{\sqrt{4-\dfrac{3}{x}+\dfrac{1}{x^2}}+a+\dfrac{b}{x}}.
		\end{eqnarray*}
		Theo giả thiết  ta có $\lim\limits_{x\to+\infty} \dfrac{\left(4-a^2\right)x-\left(3+2ab\right)+\dfrac{1-b^2}{x}}{\sqrt{4-\dfrac{3}{x}+\dfrac{1}{x^2}}+a+\dfrac{b}{x}}=0\Rightarrow \heva{&4-a^2=0\\&3+2ab=0\\&2+a\ne 0} \Leftrightarrow \heva{&a=2\\&b=-\dfrac{3}{4}.}$\\
		Do đó  $3a+8b$=0.
	} 
\end{ex}

% nguồn: Ok_Toan11_TracNghiem_Lan2
\begin{ex}%[1D3V2-5]%[ma:acb4450c]
	Giới hạn $\lim\limits_{x\to 2^+}\dfrac{2x+2}{x-2}$ có kết quả nào  sau đây?
	\choice 
		 {$-\infty$}
		{\True $+\infty$} 
		{$-\dfrac{1}{2}$} 
		{$2$} 
	 \loigiai{
		Ta  có
		\begin{align*}
			\heva{&\lim\limits_{x\to 2^+} (2x+2)=6 > 0\\&\lim\limits_{x\to 2^+} (x-2)=0\\& x-2 > 0,\forall x  > 2}
			\Rightarrow\lim\limits_{x\to 2^+}\dfrac{2x+2}{x-2}=+\infty. 
		\end{align*}
		 Vậy $\lim\limits_{x\to 2^+}\dfrac{2x+2}{x-2}=+\infty$.
	} 
\end{ex}

% nguồn: Ok_Toan11_TracNghiem_Lan2
\begin{ex}%[1D3V2-5]%[ma:6ed1245a]
	 Tìm $\lim\limits_{x\to+\infty}\left(x+1-\sqrt[3]{x^3+2}\right)$.
	 \choice
		{$-1$} 
		 {$-\infty$}
		 {$+\infty$}
		{\True $1$} 
	\loigiai{ 
	Ta có $\begin{aligned}
	 \lim\limits_{x\to+\infty}\left(1+x-\sqrt[3]{x^3+2}\right)&=\lim\limits_{x\to+\infty}\left[1+\dfrac{-2}{x^2+x\sqrt[3]{x^3+2}+\left(\sqrt[3]{x^3+2}\right)^2}\right]\\&=\lim\limits_{x\to+\infty}\left[1+\dfrac{-2}{x^2\left[1+\sqrt[3]{1+\dfrac{2}{x^3}}+\left(\sqrt[3]{1+\dfrac{2}{x^3}}\right)^2\right]}\right]\\&=\lim\limits_{x\to+\infty}\left[1+\dfrac{\dfrac{-2}{x^2}}{1+\sqrt[3]{1+\dfrac{2}{x^3}}+\left(\sqrt[3]{1+\dfrac{2}{x^3}}\right)^2}\right]=1.
\end{aligned}$\\
 Vậy $\lim\limits_{x\to+\infty}\left(x+1-\sqrt[3]{x^3+2}\right)=1$.
} 
\end{ex}

% nguồn: Ok_Toan11_TracNghiem_Lan2
\begin{ex}%[1D3V2-5]%[ma:5a0ab61d]
	Kết quả của $\lim\limits_{x\to 2^+} (x-2)\sqrt{\dfrac{x}{x^2-4}}$ bằng 
	 \choice
		{$+\infty$} 
		{\True $0$} 
		{$-\infty$} 
		{$1$} 
	 \loigiai{
		 Ta có
		\allowdisplaybreaks 
		\begin{align*}
			\lim\limits_{x\to 2^+} (x-2)\sqrt{\dfrac{x}{x^2-4}}  &=\lim\limits_{x\to 2^+}\sqrt{\dfrac{(x-2) ^2 x}{(x-2)(x+2)}}\\
			&=\lim\limits_{x\to 2^+}\sqrt{\dfrac{(x-2)x}{x+2}}=0. 
		\end{align*}
		Vậy $\lim\limits_{x\to 2^+} (x-2)\sqrt{\dfrac{x}{x^2-4}}=0$. 
	} 
\end{ex}

% nguồn: Ok_Toan11_TracNghiem_Lan2
\begin{ex}%[1D3V2-5]%[ma:e9e27326]
	Giới hạn $\lim\limits_{x\to  3^-}\dfrac{3x+2}{x-3}$ bằng
	\choice 
		{\True $-\infty$} 
		 {$+\infty$}
		{$2$} 
		{$-3$} 
	 \loigiai{
		 Ta có
		\begin{align*}
			\heva{&\lim\limits_{x\to 3^-} (3x+2)=3\cdot 3+2=11 > 0\\ &\lim\limits_{x\to 3^-} (x-3)=0\\ & x-3 < 0,\forall x <3} \Rightarrow\lim\limits_{x\to 3^-}\dfrac{3x+2}{x-3}=-\infty. 
		\end{align*}
		Vậy  $\lim\limits_{x\to 3^-}\dfrac{3x+2}{x-3}=-\infty$.
	 }
\end{ex}

% nguồn: Ok_Toan11_TracNghiem_Lan2
\begin{ex}%[1D3V2-5]%[ma:2e8123f4]
	Biết rằng $\lim\limits_{x\to+\infty}\left(\sqrt{2x^2+x}-x\sqrt{2}\right)=a\sqrt{2}+b$, với $a$, $b$ là các số hữu tỉ. Tính  $S=4a+b$.
	 \choice
		 {$S=5$}
		 {\True $S=1$}
		 {$S=-1$}
		{$S=-5$} 
	 \loigiai{
		 Ta có
		 \allowdisplaybreaks
		\begin{align*}
			\lim\limits_{x\to+\infty}\left(\sqrt{2x^2+x}-x\sqrt{2}\right) 
			 &=\lim\limits_{x\to+\infty}\dfrac{2x^2+x-2x^2}{\sqrt{2x^2+x}+x\sqrt{2}}\\
			&=\lim\limits_{x\to+\infty} \dfrac{x}{\sqrt{2x^2+x}+x\sqrt{2}}  \\
			&=\lim\limits_{x\to+\infty}\dfrac{x}{x\sqrt{2+\dfrac{1}{x}}+x\sqrt{2}} \\ 
			&=\lim\limits_{x\to+\infty}\dfrac{1}{\sqrt{2+\dfrac{1}{x}}+\sqrt{2}} \\ 
			 &=\dfrac{1}{2\sqrt{2}}=\dfrac{\sqrt{2}}{4}.
		\end{align*}
		Vậy $a=\dfrac{1}{4}$ và $b=0$.  Do đó, $S=4\cdot\dfrac{1}{4}+0=1$.
	 }
\end{ex}

% nguồn: Ok_Toan11_TracNghiem_Lan2
\begin{ex}%[1D3V2-5]%[ma:ae526448]
	Cho các số thực $a$, $b$, $c$ thỏa  mãn $c^2+a=18$ và $\lim\limits_{x\to+\infty} \left(\sqrt{ax^2+bx}-cx\right)=-2$. Tính $P=a+b+5c$.
	\choice 
		 {$P=18$}
		{\True $P=12$} 
		 {$P=9$}
		{$P=5$} 
	\loigiai{ 
		 Vì $\lim\limits_{x\to+\infty} \left(\sqrt{ax^2+bx}-cx\right)=-2\Leftrightarrow \heva{&18-2c^2=0\\&\dfrac{b}{\sqrt{a}+c}=-2.}$\\
		Mà  $c^2+a=18$.\\
		Khi đó ta có $\heva{&18-2c^2=0\\&\dfrac{b}{\sqrt{a}+c}=-2\\&c^2+a=18} \Leftrightarrow \heva{&a=9\\&c=3\\&b=-12}\Rightarrow  P=a+b+5c=9-12+5\cdot 3=12$.
	 }
\end{ex}

% nguồn: Ok_Toan11_TracNghiem_Lan2
\begin{ex}%[1D3V2-5]%[ma:9ca75254]
	Kết  quả của $\lim\limits_{x\to+\infty} (1-2x)\sqrt{\dfrac{3x-11}{x^3-1}}$ bằng
	 \choice
		 {\True $-2\sqrt{3}$}
		{$2\sqrt{3}$} 
		{$-\infty$} 
		{$0$} 
	 \loigiai{
		Khi $x\to+\infty$ thì $1-2x < 0$ nên ta  có
		 \allowdisplaybreaks
		\begin{align*}
			\lim\limits_{x\to+\infty} (1-2x)\sqrt{\dfrac{3x-11}{x^3-1}} &=\lim\limits_{x\to+\infty}\left(-\sqrt{\dfrac{(1-2x)^2 (3x-11)}{x^3-1}}\right)\\ 
			&=\lim\limits_{x\to+\infty}\left(-\sqrt{\dfrac{\left(\dfrac{1}{x}-2\right)^2\left(3-\dfrac{11}{x}\right)}{1-\dfrac{1}{x^3}}}\right)\\ 
			 &=-\sqrt{12}=-2\sqrt{3}.
		\end{align*}
	 Vậy $\lim\limits_{x\to+\infty} (1-2x)\sqrt{\dfrac{3x-11}{x^3-1}}=-2\sqrt{3}$}.
\end{ex}

% nguồn: Ok_Toan11_TracNghiem_Lan2
\begin{ex}%[1D3V2-5]%[ma:02a73b99]
	Giá  trị của giới hạn $\lim\limits_{x\to+\infty} \left(\sqrt[{3}]{2x-1}-\sqrt[{3}]{2x+1} \right)$ là
	\choice 
		 {\True $0$}
		 {$-1$}
		{$+\infty$} 
		{$-\infty$} 
	 \loigiai{
		Ta có  $\lim\limits_{x\to+\infty} \left(\sqrt[{3}]{2x-1}-\sqrt[{3}]{2x+1} \right)=\lim\limits_{x\to+\infty} \dfrac{-2}{\sqrt[{3}]{\left(2x-1\right)^2}+\sqrt[{3}]{\left(2x-1\right)\left(2x+1\right)}+\sqrt[{3}]{\left(2x+1\right)^2}}=0$.
	 }
\end{ex}

% nguồn: Ok_Toan11_TracNghiem_Lan2
\begin{ex}%[1D3V2-5]%[ma:32c9936b]
	Biết rằng $\lim\limits_{x\to-\infty}\left(\sqrt{5x^2-3x+1}+x\sqrt{5}\right)=\dfrac{a}{b}\sqrt{5}$, ($a$ là số nguyên, $b$ 
	là số nguyên dương, $\dfrac{a}{b}$  tối giản). Tổng $a+b$ có giá trị là
	 \choice
		 {$1$}
		{$5$} 
		{$4$} 
		 {\True $13$}
	\loigiai{ 
		Ta có 
		\allowdisplaybreaks 
		\begin{align*}
			\lim\limits_{x\to-\infty}\left(\sqrt{5x^2-3x+1}+x\sqrt{5}\right) 
			&=\lim\limits_{x\to-\infty}\dfrac{\left(\sqrt{5x^2-3x+1}+x\sqrt{5}\right)\left(\sqrt{5x^2-3x+1}-x\sqrt{5}\right)}{\sqrt{5x^2-3x+1}-x\sqrt{5}}\\ 
			&=\lim\limits_{x\to-\infty}\dfrac{-3x+1}{\sqrt{5x^2-3x+1}-x\sqrt{5}}\\ 
			 &=\lim\limits_{x\to-\infty}\dfrac{x\left(-3+\dfrac{1}{x}\right)}{x\left(-\sqrt{5-\dfrac{3}{x}+\dfrac{1}{x^2}}-\sqrt{5}\right)}\\
			&=\lim\limits_{x\to-\infty}\dfrac{-3+\dfrac{1}{x}}{-\sqrt{5-\dfrac{3}{x}+\dfrac{1}{x^2}}-\sqrt{5}}\\ 
			&=\dfrac{-3}{-\sqrt{5}-\sqrt{5}}=\dfrac{3}{10}\sqrt{5}. 
		\end{align*}
		 Vậy $a=3$, $b=10$. Suy ra $a+b=3+10=13$.
	 }
\end{ex}

% nguồn: Ok_Toan11_TracNghiem_Lan2
\begin{ex}%[1D3V2-5]%[ma:b52d254f]
	Biết  rằng $\lim\limits_{x\to-\infty} \left(\sqrt{2x^2-3x+1}+x\sqrt{2} \right)=\dfrac{a}{b} \sqrt{2}$, ($a,b\in \mathbb{Z},\dfrac{a}{b}$ tối giản). Tổng $a+b$ có giá trị là
	 \choice
		 {$1$}
		 {$5$}
		 {$4$}
		{\True $7$} 
	 \loigiai{
		 Ta có
		\begin{eqnarray*}
			\lim\limits_{x\to-\infty} \left(\sqrt{2x^2-3x+1}+x\sqrt{2} \right)&=&\lim\limits_{x\to-\infty} \dfrac{-3x+1}{\sqrt{2x^2-3x+1}-x\sqrt{2}}\\ &=&\lim\limits_{x\to-\infty}  \dfrac{-3+\dfrac{1}{x}}{-\sqrt{2-\dfrac{3}{x}+\dfrac{1}{x^2}}-\sqrt{2}}=\dfrac{3}{2\sqrt{2}}\\
			&=&\dfrac{3}{4}  \sqrt{2}.
		\end{eqnarray*}
		Khi đó $a=3$, $b=4$. Vậy $a+b=7$. 
	 }
\end{ex}

% nguồn: Ok_Toan11_TracNghiem_Lan2
\begin{ex}%[1D3V2-5]%[ma:379f1973]
	Giới hạn $\lim\limits_{x\to  2^+}\dfrac{3x+1}{x-2}$ bằng
	\choice 
		{\True $+\infty$} 
		{$-\infty$} 
		{$3$} 
		 {$7$}
	\loigiai{ 
		Ta  có
		\begin{align*}
			\heva{&\lim\limits_{x\to 2^+}(3x+1)=7 >0\\ &\lim\limits_{x\to 2^+}(x-2)=0\\ &  x-2 >0,\forall x >2}\Rightarrow\lim\limits_{x\to 2^+}\dfrac{3x+1}{x-2}=+\infty.
		\end{align*}
		Vậy $\lim\limits_{x\to 2^+}\dfrac{3x+1}{x-2}=+\infty$. 
	} 
\end{ex}

% nguồn: Ok_Toan11_TracNghiem_Lan2
\begin{ex}%[1D3V2-5]%[ma:4da4315c]
	 Giới hạn $\lim\limits_{x\to 1^-}\dfrac{-3x-1}{x+1}$ bằng
	\choice 
		 {$+\infty$}
		{$-\infty$} 
		{\True  $-2$}
		 {$2$}
	 \loigiai{
		Ta  có
		\begin{align*}
			 \lim\limits_{x\to 1^-}\dfrac{-3x-1}{x+1}=\dfrac{-3.1-1}{1+1}=-2.
		\end{align*}
		Vậy  $\lim\limits_{x\to 1^-}\dfrac{-3x-1}{x+1}=-2$.
	 }
\end{ex}

% nguồn: Ok_Toan11_TracNghiem_Lan2
\begin{ex}%[1D3V2-5]%[ma:1b4a56fc]
	 Giá trị của tham số $a\in \mathbb{R}$ để giới hạn $\lim\limits_{x\to+\infty} \left(\sqrt{4x^2+ax-3}-2x+1\right)=1$ là
	 \choice
		 {$a=2$}
		{$a=-2$} 
		{\True  $a=0$}
		 {$a=\pm 2$}
	 \loigiai{
		 Ta có
		\begin{eqnarray*}
			\lim\limits_{x\to+\infty} \left(\sqrt{4x^2+ax-3}-2x+1\right)&=&\lim\limits_{x\to+\infty}  \dfrac{\left(a+4\right)x-4}{\sqrt{4x^2+ax-3}+2x-1}\\
			 &=&\lim\limits_{x\to+\infty} \dfrac{\left(a+4\right)-\dfrac{4}{x}}{\sqrt{4+\dfrac{a}{x}-\dfrac{3}{x^2}}+2-\dfrac{1}{x}}\\
			&=&\dfrac{a+4}{4}. 
		\end{eqnarray*}
		Mà  $\lim\limits_{x\to+\infty} \left(\sqrt{4x^2+ax-3}-2x+1\right)=1$ nên $\dfrac{a+4}{4}=1\Leftrightarrow a=0$.
	 }
\end{ex}

% nguồn: Ok_Toan11_TracNghiem_Lan2
\begin{ex}%[1D3V2-5]%[ma:74786d46]
	Kết quả của $\lim\limits_{x\to  +\infty} \left(\sqrt{x^2+2x}-x\right)$ bằng
	 \choice
		 {$0$}
		 {$-\infty$}
		{\True $1$} 
		{$+\infty$} 
	\loigiai{ 
		Ta  có
		\begin{eqnarray*}
			\lim\limits_{x\to +\infty} \left(\sqrt{x^2+2x}-x\right)&=&  \lim\limits_{x\to +\infty} \dfrac{2x}{\sqrt{x^2+2x}+x}\\
			&=&\lim\limits_{x\to  +\infty} \dfrac{2}{\sqrt{1+\dfrac{2}{x}}+1}\\
			&=&1. 
		\end{eqnarray*}
	 }
\end{ex}

% nguồn: Ok_Toan11_TracNghiem_Lan2
\begin{ex}%[1D3V2-5]%[ma:2f73d7bf]
	Cho $I=\lim\limits_{x\to-1}(x^4-2mx+m^2+3)$. Có bao nhiêu giá trị nguyên  của tham số $m$ để $I < 12$?
	 \choice
		 {$6$}
		 {\True $5$}
		{$8$} 
		{$7$} 
	\loigiai{ 
		Ta có 
		\begin{align*}
			I &=\lim\limits_{x\to-1}(x^4-2mx+m^2+3)=(-1)^4-2m(-1)+m^2+3=m^2+2m+4.\\ 
			&I < 12\Leftrightarrow  m^2+2m+4 < 12\Leftrightarrow m^2+2m-8 <0\Leftrightarrow-4 < m <2.
		\end{align*}
		Do $m$ nguyên nên $m\in\{-3,-2,-1,0,1\}$. Vậy có tất cả $5$ giá  trị $m$ thoả mãn.
	} 
\end{ex}

% nguồn: Ok_Toan11_TracNghiem_Lan2
\begin{ex}%[1D3V2-5]%[ma:bfb30862]
	Giới  hạn $\lim\limits_{x\to a^-}\dfrac{1}{x-a}$ bằng
	\choice 
		{$-\dfrac{1}{2a}$} 
		 {$0$}
		 {$+\infty$}
		 {\True $-\infty$}
	\loigiai{ 
		Ta  có
		\begin{align*}
			\heva{&\lim\limits_{x\to a^-}1=1> 0\\ &\lim\limits_{x\to a^-}(x-a)=0\\ & x-a < 0,\forall x < a.} 
		\end{align*}
		 Vậy $\lim\limits_{x\to a^-}\dfrac{1}{x-a}=-\infty$.
	 }
\end{ex}

% nguồn: Ok_Toan11_TracNghiem_Lan2
\begin{ex}%[1D3V2-5]%[ma:1878036a]
	Tìm $\lim\limits_{x\to-\infty}  \left(\sqrt{x^2+1}+x-1\right)$. Kết quả là
	\choice 
		 {$-2$}
		{$1$} 
		 {$2$}
		{\True  $-1$}
	\loigiai{ 
		\begin{eqnarray*}
			 \lim\limits_{x\to-\infty} \left(\sqrt{x^2+1}+x-1\right)&=&\lim\limits_{x\to-\infty} \dfrac{2x}{\sqrt{x^2+1}-x+1}\\
			&=&\lim\limits_{x\to-\infty}  \dfrac{2}{-\sqrt{1+\dfrac{1}{x^2}}-1+\dfrac{1}{x}}\\
			 &=&-1.
		\end{eqnarray*}
		 Vậy $\lim\limits_{x\to-\infty} \left(\sqrt{x^2+1}+x-1\right)=-1$.
	 }
\end{ex}

% nguồn: Ok_Toan11_TracNghiem_Lan2
\begin{ex}%[1D3V2-5]%[ma:88b2fbea]
	Giới hạn $\lim\limits_{x\to  1^-}\dfrac{4x-3}{x-1}$ có kết quả nào sau đây?
	\choice 
		{\True  $-\infty$}
		{$+\infty$} 
		{$3$} 
		{$4$} 
	\loigiai{ 
		 Ta có
		\begin{align*}
			\heva{&\lim\limits_{x\to 1^-} (4x-3)=1 > 0\\ 
			 &\lim\limits_{x\to 1^-} (x-1)=0\\
			&x-1 < 0,\forall x <1} \Rightarrow\lim\limits_{x\to 1^-}\dfrac{4x-3}{x-1}=-\infty. 
		\end{align*}
		Vậy $\lim\limits_{x\to  1^-}\dfrac{4x-3}{x-1}=-\infty$.
	 }
\end{ex}

% nguồn: Ok_Toan11_TracNghiem_Lan2
\begin{ex}%[1D3V2-5]%[ma:6ac6f611]
	Cho $\lim\limits_{x\to+\infty}\left(\sqrt[3]{x^3-3ax^2+1}-bx\right)=5$, khi đó 
	\choice 
		{\True  $2a+b=-9$}
		 {$a+b=-8$}
		{$a+2b=-8$} 
		 {$a-b=-8$}
	 \loigiai{
		Ta có $b=1$ vì nếu $b\ne 1$  thì
		\begin{align*}
			\displaystyle\lim\limits_{x\to+\infty}\left(\sqrt[3]{x^3-3ax^2+1}-bx\right) 
			 =\displaystyle\lim\limits_{x\to+\infty}x\left(\sqrt[3]{1-\dfrac{3a}{x}+\dfrac{1}{x^3}}-b\right)=\pm\infty.
		\end{align*}
		 Khi đó
		 \allowdisplaybreaks
		\begin{align*}
			\displaystyle\lim\limits_{x\to+\infty}\left(\sqrt[3]{x^3-3ax^2+1}-x\right)  &
			=\displaystyle\lim\limits_{x\to+\infty}\dfrac{x^3-3ax^2+1-x^3}{\sqrt[3]{(x^3-3ax^2+1)^2}+x\sqrt[3]{x^3-3ax^2+1}+x^2}\\ 
			&=\lim\limits_{x\to+\infty}\dfrac{-3ax^2+1}{\sqrt[3]{(x^3-3ax^2+1)^2}+x\sqrt[3]{x^3-3ax^2+1}+x^2}\\ 
			 &=\lim\limits_{x\to+\infty}\dfrac{-3a+\dfrac{1}{x^2}}{\sqrt[3]{\left(1-\dfrac{3a}{x}+\dfrac{1}{x^3}\right)^2}+\sqrt[3]{1-\dfrac{3a}{x}+\dfrac{1}{x^3}}+1}\\
			&=\dfrac{-3a}{3}=5. 
		\end{align*}
		Nên $a=-5$.  Vậy $2a+b=-9$.
	} 
\end{ex}

% nguồn: Ok_Toan11_TracNghiem_Lan2
\begin{ex}%[1D3V2-5]%[ma:63c746e2]
	Kết quả của $\lim\limits_{x\to  (-1)^+}\left(x^3+1\right)\cdot\sqrt{\dfrac{x}{x^2-1}}$ bằng
	 \choice
		{$2$} 
		 {\True $0$}
		{$-\infty$} 
		{$+\infty$} 
	 \loigiai{
		Ta  có
		\begin{align*}
			\lim\limits_{x\to (-1)^+}\left(x^3+1\right)\cdot\sqrt{\dfrac{x}{x^2-1}} &=\lim\limits_{x\to (-1)^+}  (x+1)(x^2-x+1)\sqrt{\dfrac{x}{(x-1)(x+1)}}\\
			&=\lim\limits_{x\to  (-1)^+}\sqrt{x+1} (x^2-x+1)\sqrt{\dfrac{x}{x-1}}=0.
		\end{align*}
		Vậy $\lim\limits_{x\to  (-1)^+}\left(x^3+1\right)\sqrt{\dfrac{x}{x^2-1}}=0$.
	} 
\end{ex}

% nguồn: Ok_Toan11_TracNghiem_Lan2
\begin{ex}%[1D3V2-5]%[ma:07010385]
	Tính  $N=\lim\limits_{x\to 3^+}\dfrac{x+2}{x-3}$.
	\choice 
		{$-\infty$} 
		{\True  $+\infty$}
		 {$2$}
		 {$-3$}
	\loigiai{ 
		 Ta có
		\begin{align*}
			 \heva{
			 &\lim\limits_{x\to 3^+} (x+2)=3+2=5 > 0\\
			&\lim\limits_{x\to 3^+} (x-3)=0\\ 
			& x-3 >  0,\forall x >3
			} 
			\Rightarrow N=\lim\limits_{x\to 3^+}\dfrac{x+2}{x-3}=+\infty. 
		\end{align*}
		Vậy $N=\lim\limits_{x\to 3^+}\dfrac{x+2}{x-3}=+\infty$. 
	} 
\end{ex}

% nguồn: Ok_Toan11_TracNghiem_Lan2
\begin{ex}%[1D3V2-5]%[ma:0af8dd29]
	Biết rằng tồn tại duy nhất một số thực $m$ để $\lim\limits_{x\to-\infty} \left(\sqrt{x^2+mx+1}+x\right)=-\dfrac{1}{4}$. Mệnh đề nào sau đây đúng? 
	 \choice
		{$m\in \left(1;2\right)$} 
		 {$m\in \left(-1;0\right)$}
		{\True  $m\in \left(0;1\right)$}
		{$m\in  \left(-2;-1\right)$}
	\loigiai{ 
		Ta có \begin{eqnarray*}
		 \lim\limits_{x\to-\infty} \left(\sqrt{x^2+mx+1}+x\right)&=&\lim\limits_{x\to-\infty} \dfrac{mx+1}{\sqrt{x^2+mx+1}-x}\\
		&=&\lim\limits_{x\to-\infty}  \dfrac{m+\dfrac{1}{x}}{-\sqrt{1+\dfrac{m}{x}+\dfrac{1}{x^2}}-1}\\
		 &=&-\dfrac{m}{2}.
	\end{eqnarray*}
	 Nên $m=\dfrac{1}{2} \in (0;1)$.
 }
\end{ex}

% nguồn: Ok_Toan11_TracNghiem_Lan2
\begin{ex}%[1D3V2-5]%[ma:44cf0027]
	Cho $\lim\limits_{x\to-\infty} \left(\sqrt{x^2+2ax+1}+x\right)=1$ thì giá trị của $a$ thuộc khoảng nào trong các  khoảng sau?
	\choice 
		 {$\left(-5;-2\right]$}
		{$\left[1;3\right)$} 
		 {$\left[3;5\right)$}
		 {\True $\left(-2;1\right)$}
	 \loigiai{
		Ta có $\lim\limits_{x\to-\infty}  \left(\sqrt{x^2+2ax+1}+x\right)=\lim\limits_{x\to-\infty} \dfrac{2ax+1}{\sqrt{x^2+2ax+1}-x}=-a$.\\
		Mà $\lim\limits_{x\to-\infty} \left(\sqrt{x^2+2ax+1}+x\right)=1$  nên $a=-1$.
	 }
\end{ex}

% nguồn: Ok_Toan11_TracNghiem_Lan2
\begin{ex}%[1D3V2-5]%[ma:93ad49ec]
	 Giới hạn $\lim\limits_{x\to 2^+}\dfrac{x^2+x-1}{x^2-4}$ bằng
	 \choice
		{$-\infty$} 
		{$3$} 
		 {$0$}
		{\True  $+\infty$}
	 \loigiai{
		Ta  có
		\begin{align*}
			\heva{&\lim\limits_{x\to 2^+}(x^2+x-1)=5 >0\\  &\lim\limits_{x\to 2^+}(x^2-4)=0\\ & x^2-4 >0,\forall x >2}\Rightarrow\lim\limits_{x\to 2^+}\dfrac{x^2+x-1}{x^2-4}=+\infty.
		\end{align*}
		 Vậy $\lim\limits_{x\to 2^+}\dfrac{x^2+x-1}{x^2-4}=+\infty.$
	 }
\end{ex}

% nguồn: Ok_Toan11_TracNghiem_Lan2
\begin{ex}%[1D3V2-5]%[ma:9b75bf48]
	 Tính $\lim\limits_{x\to-\infty}\dfrac{1}{x}\left(\sqrt{4x^2+1}-x\right)$.
	\choice 
		 {$+\infty$}
		 {$0$}
		 {$\dfrac{1}{2}$}
		{\True  $-3$}
	\loigiai{ 
		Ta  có
		 \allowdisplaybreaks
		\begin{align*}
			 \lim\limits_{x\to-\infty}\dfrac{1}{x}\left(\sqrt{4x^2+1}-x\right) &=\lim\limits_{x\to-\infty}\dfrac{1}{x}\left(\dfrac{4x^2+1-x^2}{\sqrt{4x^2+1}+x}\right)\\
			 &=\lim\limits_{x\to-\infty}\dfrac{3x^2+1}{x\sqrt{4x^2+1}+x^2}\\
			 &=\lim\limits_{x\to-\infty}\dfrac{3+\dfrac{1}{x^2}}{-\sqrt{4+\dfrac{1}{x^2}}+1}=-3.
		\end{align*}
		 Vậy $\lim\limits_{x\to-\infty}\dfrac{1}{x}\left(\sqrt{4x^2+1}-x\right)=-3$.
	 }
\end{ex}

% nguồn: Ok_Toan11_TracNghiem_Lan2
\begin{ex}%[1D3V2-5]%[ma:181aeb0e]
	$\lim\limits_{x\to-\infty} \left(\sqrt{4x^2-6x+5}+2x\right)$  bằng
	\choice 
		 {$\dfrac{1}{2}$}
		{$0$} 
		 {$-\dfrac{3}{2}$}
		{\True  $\dfrac{3}{2}$}
	 \loigiai{
		Ta  có
		\begin{eqnarray*}
			 \lim\limits_{x\to-\infty} \left(\sqrt{4x^2-6x+5}+2x\right)&=&\lim\limits_{x\to-\infty} \dfrac{-6x+5}{\sqrt{4x^2-6x+5}-2x}\\
			&=&\lim\limits_{x\to-\infty} \dfrac{-6+\dfrac{5}{x}}{-\sqrt{4-\dfrac{6}{x}+\dfrac{5}{x^2}}-2}\\ 
			 &=&\dfrac{3}{2}.
		\end{eqnarray*}
		Vậy  $\lim\limits_{x\to-\infty} \left(\sqrt{4x^2-6x+5}+2x\right)=\dfrac{3}{2}$.
	} 
\end{ex}

% nguồn: Ok_Toan11_TracNghiem_Lan2
\begin{ex}%[1D3V2-5]%[ma:4bf1a703]
	Tính $\lim\limits_{x\to 1^+} (x^3-1)\sqrt{\dfrac{x}{x^2-1}}$. 
	\choice 
		 {$+\infty$}
		{\True $0$} 
		{$-\infty$} 
		 {$\dfrac{3}{2}$}
	\loigiai{ 
		 Ta có
		 \allowdisplaybreaks
		\begin{align*}
			\lim\limits_{x\to 1^+} (x^3-1)\sqrt{\dfrac{x}{x^2-1}} &=\lim\limits_{x\to  1^+}\sqrt{\dfrac{(x^3-1)^2 x}{x^2-1}}\\
			 &=\lim\limits_{x\to 1^+}\sqrt{\dfrac{(x-1)^2 (x^2+x+1)^2 x}{(x-1)(x+1)}}\\
			 &=\lim\limits_{x\to 1^+}\sqrt{\dfrac{(x-1) (x^2+x+1)^2 x}{x+1}}=0.
		\end{align*}
		Vậy $\lim\limits_{x\to 1^+} (x^3-1)\sqrt{\dfrac{x}{x^2-1}}=0$. 
	} 
\end{ex}

% nguồn: Ok_Toan11_TracNghiem_Lan2
\begin{ex}%[1D3V2-5]%[ma:a83ea8ad]
	$\lim\limits_{x\to+\infty} \left(\sqrt{x+1}-\sqrt{x-3}  \right)$ bằng
	 \choice
		 {\True $0$}
		{$2$} 
		{$-\infty$} 
		{$+\infty$} 
	\loigiai{ 
		\begin{eqnarray*}
			\lim\limits_{x\to+\infty} \left(\sqrt{x+1}-\sqrt{x-3} \right)&=&\lim\limits_{x\to+\infty} \dfrac{x+1-x+3}{\sqrt{x+1}+\sqrt{x-3}}\\ 
			&=&\lim\limits_{x\to+\infty}  \dfrac{4}{\sqrt{x+1}+\sqrt{x-3}}\\
			&=&0. 
		\end{eqnarray*}
		 Vậy $\lim\limits_{x\to+\infty} \left(\sqrt{x+1}-\sqrt{x-3} \right)=0$.
	 }
\end{ex}

% nguồn: Ok_Toan11_TracNghiem_Lan2
\begin{ex}%[1D3V2-5]%[ma:ee948a7e]
	Trong các giới hạn sau,  giới hạn nào bằng $0$?
	\choice 
		 {$\lim\limits_{x\to 1} \dfrac{x-1}{x^3-1}$}
		{$\lim\limits_{x\to  -2} \dfrac{2x+5}{x+10}$}
		{$\lim\limits_{x\to -2} \dfrac{x^2-1}{x^2+3x-2}$} 
		{\True  $\lim\limits_{x\to+\infty} \left(\sqrt{x^2+1}-x\right)$}
	 \loigiai{
		\begin{itemize}
			 \item $\lim\limits_{x\to 1} \dfrac{x-1}{x^3-1}=\lim\limits_{x\to 1} \dfrac{x-1}{\left(x-1\right)\left(x^2+x+1\right)}=\lim\limits_{x\to 1} \dfrac{1}{x^2+x+1}=\dfrac{1}{3}$.
			\item  $\lim\limits_{x\to -2} \dfrac{2x+5}{x+10}=\dfrac{1}{8}$.
			 \item $\lim\limits_{x\to -2} \dfrac{x^2-1}{x^2+3x-2}=-\dfrac{3}{4}$.
			\item  $\lim\limits_{x\to+\infty} \left(\sqrt{x^2+1}-x\right)=\lim\limits_{x\to+\infty} \dfrac{1}{\sqrt{x^2+1}+x}=0$.
		\end{itemize}
		Vậy $\lim\limits_{x\to+\infty} \left(\sqrt{x^2+1}-x\right)=0$. 
	} 
\end{ex}

% nguồn: Ok_Toan11_TracNghiem_Lan2
\begin{ex}%[1D3V2-5]%[ma:19e7bd9a]
	Giới  hạn $\lim\limits_{x\to 1^+}\dfrac{2x-1}{x-1}$ bằng
	\choice 
		{\True $+\infty$} 
		 {$1$}
		 {$0$}
		{$-\infty$} 
	\loigiai{ 
		Ta có 
		\begin{align*}
			\heva{&\lim\limits_{x\to 1^+} (2x-1)=1 > 0\\ & \lim\limits_{x\to 1^-} (x-1)=0\\& x-1 > 0,\forall x >1}\Rightarrow\lim\limits_{x\to 1^+}\dfrac{2x-1}{x-1}=+\infty. 
		\end{align*}
		Vậy  $\lim\limits_{x\to 1^+}\dfrac{2x-1}{x-1}=+\infty$.
	 }
\end{ex}

% nguồn: Ok_Toan11_TracNghiem_Lan2
\begin{ex}%[1D3V2-5]%[ma:d49cc3d0]
	Biết rằng $\lim\limits_{x\to+\infty}  \left(\sqrt{5x^2+2x}-x\sqrt{5} \right)=a\sqrt{5}+b$. Tính $S=5a+b$.
	 \choice
		 {$S=5$}
		{\True $S=1$} 
		{$S=-1$} 
		{$S=-5$} 
	\loigiai{ 
		Ta có 
		\begin{eqnarray*}
			 \lim\limits_{x\to+\infty} \left(\sqrt{5x^2+2x}-x\sqrt{5} \right)&=&\lim\limits_{x\to+\infty} \dfrac{2x}{\sqrt{5x^2+2x}+x\sqrt{5}}\\
			&=&\lim\limits_{x\to+\infty}  \dfrac{2}{\sqrt{5+\dfrac{2}{x}}+\sqrt{5}}\\
			&=&\dfrac{1}{5} \sqrt{5}. 
		\end{eqnarray*}
		Theo giả thiết $\lim\limits_{x\to+\infty} \left(\sqrt{5x+2x}-x\sqrt{5}  \right)=a\sqrt{5}+b$. Suy ra $\heva{&a=\dfrac{1}{5}\\&b=0.}$ \\
		Do đó $S=5a+b=1$. 
	} 
\end{ex}

% nguồn: Ok_TracNghiemLop11_Lan1
\begin{ex}%[1D3V2-5]%[ma:4fe1084e]
	Giá trị  của giới hạn $\lim\limits_{x\to+\infty}\left(\sqrt{x^2+x}-\sqrt[3]{x^3-x^2}\right)$ là
	\choice 
		{\True  $\dfrac{5}{6}$}
		 {$+\infty $}
		 {$-1$}
		{$-\infty  $}
	 \loigiai{
		\begin{eqnarray*}
			\lim\limits_{x\to+\infty}\left(\sqrt{x^2+x}-\sqrt[3]{x^3-x^2}\right)&=&  \lim\limits_{x\to+\infty}\left(\sqrt{x^2+x}-x+x-\sqrt[3]{x^3-x^2}\right)\\
			&=&\lim\limits_{x\to+\infty}\left(\dfrac{x}{\sqrt{x^2+1}+x}+\dfrac{x^2}{x^2+x\sqrt[3]{x^3-1}+\sqrt[3]{\left(x^3-1\right)^2}}\right)\\ &=&\dfrac{1}{2}+\dfrac{1}{3}=\dfrac{5}{6}. 
		\end{eqnarray*}
	} 
\end{ex}

% nguồn: Ok_TracNghiemLop11_Lan1
\begin{ex}%[1D3B2-5]%[ma:8c901b5d]
	Tính giới hạn $L=\displaystyle\lim\limits_{x\to 1}  \dfrac{\sqrt{x+3}-2}{x-1}$.
	 \choice
		{\True $L=\dfrac{1}{4}$} 
		 {$L=1$}
		{$L=2$} 
		{$L=\dfrac{1}{2}$} 
	\loigiai{ 
		Ta có $\displaystyle\lim\limits_{x\to  1} \dfrac{\sqrt{x+3}-2}{x-1}=\displaystyle\lim\limits_{x\to 1}\dfrac{x-1}{(x-1)\left(\sqrt{x+3}+2\right)}=\displaystyle\lim\limits_{x\to 1}\dfrac{1}{\sqrt{x+3}+2}=\dfrac{1}{4}$.
	} 
\end{ex}


%%%%%%%%%%%%%%%% ĐÚNG/SAI %%%%%%%%%%%%%%%%

% nguồn: Ok_dungsai_lop11_Lan1
\begin{ex}%[1D3N2-5]%[ma:c29d146a]
 Mỗi kết quả giới hạn sau đúng  hay sai?
 \choiceTF 
  {\True $\lim\limits_{x \to +\infty} \left(x-\sqrt{x^2+2}\right)=0$} 
  {\True  $\lim\limits_{x \to +\infty} \left(\sqrt{x^2 + 2x}-x\right)=1$}
  {$\lim\limits_{x \to +\infty}  \left(\sqrt{2x^2 + 2x}-x\right)=2$}
  {$\lim\limits_{x \to -\infty} \left(\sqrt{x^2 +  2x}+x\right)=2$}
  \loigiai{
  \begin{itemchoice}
    \itemch $\lim\limits_{x \to +\infty} \left(x-\sqrt{x^2+2}\right)
   =\lim\limits_{x  \to +\infty} \dfrac{-2}{x+\sqrt{x^2+2}}
   =\lim\limits_{x  \to +\infty} \dfrac{-2}{x\left(1+\sqrt{1+\dfrac{2}{x^2}}\right)}
   =0$. 
   \itemch $\lim\limits_{x  \to +\infty} \left(\sqrt{x^2 + 2x}-x\right)
   =\lim\limits_{x \to  +\infty} \dfrac{2x}{\sqrt{x^2+2x}+x}
   =\lim\limits_{x \to +\infty} \dfrac{2}{\sqrt{1+\dfrac{2}{x}}+1} 
    =1$.
   \itemch $\lim\limits_{x \to +\infty}  \left(\sqrt{2x^2 + 2x}-x\right)=\lim\limits_{x \to +\infty} x\left(\sqrt{2+\dfrac{2}{x}}-1\right)=+\infty$, vì $\lim\limits_{x \to +\infty} \left(\sqrt{2+\dfrac{2}{x}}-1\right)=\sqrt{2}-1>0$.
    \itemch
    \allowdisplaybreaks
   \begin{eqnarray*}
    \lim\limits_{x  \to -\infty} \left(\sqrt{x^2 + 2x}+x\right)
     &=&\lim\limits_{x \to -\infty} \dfrac{2x}{\sqrt{x^2+2x}-x}
     =\lim\limits_{x \to -\infty} \dfrac{2x}{-x\sqrt{1+\dfrac{2}{x}}-x}\\
    &=&\lim\limits_{x \to -\infty} \dfrac{2}{-\sqrt{1+\dfrac{2}{x}}-1}=-1. 
   \end{eqnarray*}
  \end{itemchoice}
 } 
\end{ex}

% nguồn: Ok_dungsai_lop11_Lan1
\begin{ex}%[1D3H2-5]%[ma:4106edde]
 Mỗi kết quả giới hạn sau đúng hay sai? 
 \choiceTF 
  {$\lim\limits_{x \to 1} \dfrac{\sqrt{2x + 1}-\sqrt{x +  2}}{x-1}=1$}
  {$\lim\limits_{x\to  2} \dfrac{\sqrt{x+2}-2}{x-2}=-\dfrac{1}{4}$}
  {$\lim\limits_{x \to 2}\dfrac{2 - x}{\sqrt{x + 7}  - 3}=6$}
  {\True $\lim\limits_{x \to 0} \dfrac{x}{\sqrt{x+1}-1}=2$} 
 \loigiai{ 
  \begin{itemchoice}
   \itemch  $\lim\limits_{x \to 1} \dfrac{\sqrt{2x+1}-\sqrt{x+2}}{x-1}=\lim\limits_{x \to 1}\dfrac{x-1}{(x-1)\left(\sqrt{2x+1}-\sqrt{x+2} \right)}=\lim\limits_{x \to 1}\dfrac{1}{\sqrt{2x+1}+\sqrt{x+2}}=\dfrac{1}{2\sqrt{3}}$.
   \itemch $\lim\limits_{x\to  2}\dfrac{\sqrt{x+2}-2}{x-2}=\lim\limits_{x\to 2}\dfrac{x-2}{(x-2)\left(\sqrt{x+2}+2\right)}=\lim\limits_{x\to 2}\dfrac{1}{\sqrt{x+2}+2}=\dfrac{1}{4}$.
   \itemch $\lim\limits_{x \to 2}\dfrac{2 -  x}{\sqrt{x + 7} - 3}
   =\lim\limits_{x \to  2}\dfrac{(2 - x)\left(\sqrt{x + 7} + 3\right)}{x + 7 - 9}
   = \lim\limits_{x \to 2}\dfrac{(2  - x)\left(\sqrt{x + 7} + 3\right)}{- (2 - x)}= - 6$.
   \itemch $\lim\limits_{x  \to 0} \dfrac{x}{\sqrt{x+1}-1}
   =\lim\limits_{x \to 0}  \left(\sqrt{x+1}+1\right)=2$.
  \end{itemchoice}
 } 
\end{ex}

% nguồn: Ok_dungsai_lop11_Lan1
\begin{ex}%[1D3H2-5]%[ma:8add3723]
 Xác định tính đúng, sai của các  mệnh đề sau:
 \choiceTF 
  {\True Kết quả của $\lim \limits_{x \rightarrow+\infty}  \dfrac{x-1}{2 x+3}$ là $\dfrac{1}{2}$}
  {\True Kết quả của $\lim \limits_{x \rightarrow+\infty} \dfrac{\sqrt{x}+1}{x-1}$ là  0}
  {Kết  quả của $\lim \limits_{x \rightarrow-\infty}\left(\sqrt{x^2-3 x+7}+x\right)$ là 2}
   {\True Kết quả của $\lim \limits_{x \rightarrow-\infty}\left(\sqrt{x^2+3 x-1}+\sqrt[3]{x^3-2 x^2+2}\right)$ là $\dfrac{-13}{6}$}
 \loigiai{ 
  \begin{itemchoice}
   \itemch Đúng. Ta có $\lim  \limits_{x \rightarrow+\infty} \dfrac{x-1}{2 x+3}=\lim \limits_{x \rightarrow+\infty} \dfrac{1-\dfrac{1}{x}}{2+\dfrac{3}{x}}=\dfrac{1}{2}$.
   \itemch Đúng. Ta có $\lim \limits_{x  \rightarrow+\infty} \dfrac{\sqrt{x}+1}{x-1}=\lim \limits_{x \rightarrow+\infty} \dfrac{\dfrac{1}{\sqrt{x}}+\dfrac{1}{x}}{1-\dfrac{1}{x}}=0$.
   \itemch Sai.  Ta có
   \begin{align*}
    \lim \limits_{x \rightarrow-\infty}\left(\sqrt{x^2-3 x+7}+x\right)=\lim \limits_{x \rightarrow-\infty} \dfrac{-3 x+7}{\sqrt{x^2-3  x+7}-x}=\lim \limits_{x \rightarrow-\infty} \dfrac{-3+\dfrac{7}{x}}{-\sqrt{1-\dfrac{3}{x}+\dfrac{7}{x^2}}-1}=\dfrac{3}{2}.
   \end{align*}
   \itemch Đúng.  Ta có
   \begin{align*}
    &\lim \limits_{x \rightarrow-\infty}\left(\sqrt{x^2+3  x-1}+\sqrt[3]{x^3-2 x^2+2}\right)\\
    =&\lim \limits_{x  \rightarrow-\infty}\left(\sqrt{x^2+3 x-1}+x-x+\sqrt[3]{x^3-2 x^2+2}\right)\\
    =& \lim  \limits_{x \rightarrow-\infty}\left(\dfrac{\left(\sqrt{x^2+3 x-1}+x\right)\left(\sqrt{x^2+3 x-1}-x\right)}{\sqrt{x^2+3 x-1}-x}+\dfrac{\left(\sqrt[3]{x^3-2 x^2+2}\right)^3-x^3}{\sqrt[3]{\left(x^3-2 x^2+2\right)^2}+x \sqrt[3]{x^3-2 x^2+2}+x^2}\right)\\
    =&\lim \limits_{x \rightarrow-\infty}\left(\dfrac{3  x-1}{\sqrt{x^2+3 x-1}-x}+\dfrac{-2 x^2+2}{\sqrt[3]{\left(x^3-2 x^2+2\right)^2}+x \sqrt[3]{x^3-2 x^2+2}+x^2}\right)\\
     =&\lim \limits_{x \rightarrow-\infty}\left(\dfrac{3-\dfrac{1}{x}}{-\sqrt{1+\dfrac{3}{x}-\dfrac{1}{x^2}}-1}+\dfrac{-2+\dfrac{2}{x^2}}{\sqrt[3]{\left(1-\dfrac{2}{x}+\dfrac{2}{x^3}\right)^2}+\sqrt[3]{1-\dfrac{2}{x}+\dfrac{2}{x^3}}+1}\right)\\
    =&-\dfrac{3}{2}-\dfrac{2}{3}\\ 
    =&-\dfrac{13}{6}. 
   \end{align*}
  \end{itemchoice}
  }
\end{ex}

% nguồn: Ok_dungsai_lop11_Lan1
\begin{ex}%[1D3H2-5]%[ma:26ec056a]
  Tìm được các giới hạn sau.
  \choiceTF
   {\True $\lim\limits_{x\to 0}\dfrac{\sqrt{4+x}-2}{4x}=\dfrac{1}{16}$}
  {\True  $\lim\limits_{x\to 2}\dfrac{4-x^2}{\sqrt{x+7}-3}=-24$}
  {\True  $\lim\limits_{x\to 2}\dfrac{\sqrt{2x+5}-3}{\sqrt{x+2}-2}=\dfrac{4}{3}$}
   {$\lim\limits_{x\to 1}\dfrac{\sqrt[3]{x+7}-2}{x-1}=\dfrac{1}{3}$}
  \loigiai{
  \begin{itemchoice}
   \itemch \begin{eqnarray*}
   \lim\limits_{x\to 0}\dfrac{\sqrt{4+x}-2}{4x}&=& \lim\limits_{x\to 0}\dfrac{(\sqrt{4+x}-2)(\sqrt{4+x}+2)}{4x(\sqrt{4+x}+2)}=\lim\limits_{x\to  0}\dfrac{4+x-4}{4x(\sqrt{4+x}+2)}\\
   &=&\lim\limits_{x\to 0}\dfrac{1}{4(\sqrt{4+x}+2)}= \dfrac{1}{4(\sqrt 4+2)}=\dfrac{1}{16}. 
  \end{eqnarray*}
  \itemch \begin{eqnarray*}
   \lim\limits_{x\to 2}\dfrac{4-x^2}{\sqrt{x+7}-3}&=& \lim\limits_{x\to 2}\dfrac{(2-x)(2+x)(\sqrt{x+7}+3)}{(\sqrt{x+7}-3)(\sqrt{x+7}+3)}\\
  &=&  \lim\limits_{x\to 2}\dfrac{(2-x)(2+x)(\sqrt{x+7}+3)}{x+7-9}\\
  &=&\lim\limits_{x\to 2}[-(2+x)(\sqrt{x+7}+3)]=-4  \cdot 6=-24.
 \end{eqnarray*}
 \itemch \begin{eqnarray*}
 \lim\limits_{x\to 2}\dfrac{\sqrt{2x+5}-3}{\sqrt{x+2}-2}&=& \lim\limits_{x\to 2}\dfrac{(\sqrt{2x+5}-3)(\sqrt{2x+5}+3)(\sqrt{x+2}+2)}{(\sqrt{x+2}-2)(\sqrt{x+2}+2)(\sqrt{2x+5}+3)}\\ 
 &=& \lim\limits_{x\to 2}\dfrac{(2x+5-9)(\sqrt{x+2}+2)}{(x+2-4)(\sqrt{2x+5}+3)}\\  &=&\lim\limits_{x\to 2}\dfrac{2(\sqrt{x+2}+2)}{\sqrt{2x+5}+3}=\dfrac{4}{3}.
 \end{eqnarray*}
 \itemch \begin{eqnarray*}
 \lim\limits_{x\to 1}\dfrac{\sqrt[3]{x+7}-2}{x-1}&=&  \lim\limits_{x\to 1}\dfrac{(\sqrt[3]{x+7}-2)\left(\sqrt[3]{(x+7)^2}+2\sqrt[3]{x+7}+4\right)}{(x-1)\left(\sqrt[3]{(x+7)^2}+2\sqrt[3]{x+7}+4\right)}\\
 &=&  \lim\limits_{x\to 1}\dfrac{x+7-2^3}{(x-1)\left(\sqrt[3]{(x+7)^2}+2\sqrt[3]{x+7}+4\right)}\\
 &=&\lim\limits_{x\to  1}\dfrac{1}{\sqrt[3]{(x+7)^2}+2\sqrt[3]{x+7}+4}=\dfrac{1}{12}.
 \end{eqnarray*}
 \end{itemchoice}
} 
\end{ex}

% nguồn: Ok_dungsai_lop11_Lan1
\begin{ex}%[1D3H2-5]%[ma:0138e382]
 Các mệnh đề sau đây đúng hay sai? 
 \choiceTF 
  {$\lim  \dfrac{\sqrt{16n^2-2}}{n}=0$}
  {\True $\lim \dfrac{4^{n}-5^{n}}{16\cdot5^{n}-3^{n}+1}<0$} 
  {\True $\lim\limits_{x \rightarrow -1}{\dfrac{x^2-2x-3}{x+1}} = -4$} 
  {$\lim\limits_{x \rightarrow  1} \dfrac{x^2-5x+4}{\sqrt{3x+1}-2}=4$}
 \loigiai{ 
  \begin{itemchoice}
   \itemch  {\bf Sai.} \\
   $\lim \dfrac{\sqrt{16n^2-2}}{n}=\lim  \sqrt{\dfrac{16n^2-2}{n^2}}=\sqrt{\lim \left(16-2\cdot \dfrac{1}{n^2}\right)}=\sqrt{16-2\cdot 0}=4$.
   \itemch  {\bf Đúng.}
   \begin{eqnarray*}
    \lim  \dfrac{4^{n}-5^{n}}{16\cdot5^{n}-3^{n}+1}&=&
    \lim\dfrac{5^n\cdot \left[\left(\dfrac{4}{5}\right)^n-1\right]}{5^n\cdot \left[16-\left(\dfrac{3}{5}\right)^n+\left(\dfrac{1}{5}\right)^n\right]}\\ 
    &=& 
     \lim\dfrac{\left(\dfrac{4}{5}\right)^n-1}{ 16-\left(\dfrac{3}{5}\right)^n+\left(\dfrac{1}{5}\right)^n}=\dfrac{-1}{16}<0.
   \end{eqnarray*}
    \itemch {\bf Đúng.}\\
   $\lim\limits_{x \rightarrow -1}{\dfrac{x^2-2x-3}{x+1}} = \lim\limits_{x \rightarrow -1}{\dfrac{(x+1)(x-3)}{x+1}}  = \lim\limits_{x \rightarrow -1}{(x-3)} = -4$.
   \itemch {\bf  Sai.}
   \begin{eqnarray*}
    \lim\limits_{x \rightarrow  1} \dfrac{x^2-5x+4}{\sqrt{3x+1}-2}& =&
    \lim\limits_{x \rightarrow 1}  \dfrac{(x^2-5x+4)(\sqrt{3x+1}+2)}{(\sqrt{3x+1}-2)(\sqrt{3x+1}+2)}\\
     &=&
    =\lim\limits_{x \rightarrow  1} \dfrac{(x-1)(x-4)(\sqrt{3x+1}+2)}{3(x-1)} \\
    &=&  \lim\limits_{x \rightarrow 1} \dfrac{(x-4)(\sqrt{3x+1}+2)}{3}=
    \dfrac{(-3) \cdot 4}{3}=-4. 
   \end{eqnarray*}
  \end{itemchoice}
 } 
\end{ex}

% nguồn: Ok_dungsai_lop11_Lan1
\begin{ex}%[1D3V2-5]%[ma:b16b30f4]
 Xác định tính đúng, sai của các  mệnh đề sau:
  \choiceTF
  {\True Kết quả  của $\lim \limits_{x \rightarrow 2} \dfrac{(\sqrt{x}-2)(x+2)}{\sqrt{x}-2}$ là 4}
  {Kết quả của $\lim  \limits_{x \rightarrow 1} \dfrac{\sqrt{x}-1}{\sqrt{x+3}-2}$ là 3}
  {Kết  quả của $\lim \limits_{x \rightarrow 2} \dfrac{\sqrt{x+7}-3}{\sqrt{x-1}-1}$ là 2}
  {Kết quả của $\lim \limits_{x \rightarrow-1} \dfrac{2-\sqrt{3-x}}{\sqrt[3]{x}+1}$ là  3}
  \loigiai{
  \begin{itemchoice}
   \itemch  Đúng. Ta có $\lim \limits_{x \rightarrow 2} \dfrac{(\sqrt{x}-2)(x+2)}{\sqrt{x}-2}=\lim \limits_{x \rightarrow 2}(x+2)=4$.
   \itemch Sai. Ta  có
   \begin{align*}
    \lim \limits_{x \rightarrow 1}  \dfrac{\sqrt{x}-1}{\sqrt{x+3}-2}=&\lim \limits_{x \rightarrow 1} \dfrac{(\sqrt{x}-1)(\sqrt{x}+1)(\sqrt{x+3}+2)}{(\sqrt{x}+1)(\sqrt{x+3}-2)(\sqrt{x+3}+2)}\\
    =&\lim  \limits_{x \rightarrow 1} \dfrac{(x-1)(\sqrt{x+3}+2)}{(\sqrt{x}+1)(x-1)}\\
    =&\lim  \limits_{x \rightarrow 1} \dfrac{\sqrt{x+3}+2}{\sqrt{x}+1}=2.
   \end{align*}
   \itemch  Sai. Ta có
   \begin{align*}
    \lim \limits_{x \rightarrow 2} \dfrac{\sqrt{x+7}-3}{\sqrt{x-1}-1}=&\lim \limits_{x \rightarrow 2}  \dfrac{(\sqrt{x-1}+1)(\sqrt{x+7}-3)(\sqrt{x+7}+3)}{(\sqrt{x-1}+1)(\sqrt{x-1}-1)(\sqrt{x+7}+3)}\\
    =&\lim \limits_{x  \rightarrow 2} \dfrac{(x-2)(\sqrt{x-1}+1)}{(\sqrt{x+7}+3)(x-2)}\\
    =&\lim \limits_{x  \rightarrow 2} \dfrac{\sqrt{x-1}+1}{\sqrt{x+7}+3}=\dfrac{1}{3}.
   \end{align*}
   \itemch Sai.  Ta có
   \begin{align*}
    \lim \limits_{x \rightarrow-1} \dfrac{2-\sqrt{3-x}}{\sqrt[3]{x}+1}=\lim \limits_{x \rightarrow-1} \dfrac{(1+x)\left(\sqrt[3]{x^2}-\sqrt[3]{x}+1\right)}{(x+1)(2+\sqrt{3-x})}=\lim \limits_{x \rightarrow-1} \dfrac{\left(\sqrt[3]{x^2}-\sqrt[3]{x}+1\right)}{(2+\sqrt{3-x})}=\dfrac{3}{4}. 
   \end{align*}
  \end{itemchoice}
  }
\end{ex}

% nguồn: Ok_dungsai_lop11_Lan2
\begin{ex}%[1D3V2-5]%[ma:c3b0823e]
	Biết  $\lim\limits_{x\to+\infty }\left(\sqrt{4x^2-3x+1}-\left(ax+b\right)\right)=0$.
	 \choiceTF
		{$a$ là số  lẻ}
		{$  b>0$}
		{\True $ a\cdot b<0$} 
		{\True $  a-4b=5$}
	\loigiai{Ta  có
		\begin{align*}
			\lim\limits_{x\to+\infty }\left(\sqrt{4x^2-3x+1}-\left(ax+b\right)\right)=0&\Leftrightarrow\lim\limits_{x\to+\infty }\left(\left(\sqrt{4x^2-3x+1}-ax\right)-b\right)=0\\&\Leftrightarrow\lim\limits_{x\to+\infty }\left(\dfrac{4x^2-3x+1-{a^2}x^2}{\sqrt{4x^2-3x+1}+ax}-b\right)=0\\&\Leftrightarrow\lim\limits_{x\to+\infty  }\left(\dfrac{\left(4-{a^2}\right)x^2-3x+1}{\sqrt{4x^2-3x+1}+ax}-b\right)=0\\&\Leftrightarrow\heva{&4-{a^2}=0\\&a>0\\&\dfrac{-3}{2+a}-b=0\\
			}\\&\Leftrightarrow\heva{&a=2\\&b=-\dfrac34. 
			} 
		\end{align*}
		\begin{itemchoice}
			\itemch $a$  là một số chẵn.
			 \itemch $ b<0$.
			\itemch $  a\cdot b<0$.
			\itemch  $ a-4b=5$.
	\end{itemchoice}}
\end{ex}

% nguồn: Ok_dungsai_lop11_Lan2
\begin{ex}%[1D3V2-5]%[ma:2cc14b4b]
	Cho các số thực $a$, $b$, $c$ thỏa  mãn $c^2+a=18$ và $\lim\limits_{x\to+\infty }\left(\sqrt{ax^2+bx}-cx\right)=-2$.
	 \choiceTF
		{\True $ a=9$} 
		{$  b=3c$}
		 {\True $ a=3c$}
		{$P=a+b+5c=14$} 
	\loigiai{Nếu $c\le0$ thì $\lim\limits_{x\to+\infty }\left(\sqrt{ax^2+bx}-cx\right)=+\infty$.\\ 
		Với $c>0$, ta  có
		\begin{align*}
			\lim\limits_{x\to+\infty }\left(\sqrt{ax^2+bx}-cx\right)&=\lim\limits_{x\to+\infty  }\dfrac{\left(a-c^2\right)x^2+bx}{\sqrt{ax^2+bx}+cx}\\
			&=\lim\limits_{x\to+\infty}\dfrac{(a-c^2)x+b}{\sqrt{a+\dfrac bx}+c}. 
		\end{align*}
		 Do đó $\lim\limits_{x\to+\infty }\left(\sqrt{ax^2+bx}-cx\right)=-2$ chỉ xảy ra khi
		$\Leftrightarrow\heva{&a-c^2=0\\&c>0\\&\dfrac b{\sqrt a+c}=-2.}$\\ 
		Mặt  khác, ta cũng có $c^2+a=18$.\\
		Do  đó, $\heva{&a=c^2=9\\&c>0\\&b=-2\left(\sqrt a+c\right)}\Leftrightarrow\heva{&a=9\\&b=-12\\&c=3.}$\\
		 Vậy $P=a+b+5c=12$.
		\begin{itemchoice}
			\itemch  $ a=9$.
			\itemch $b=-4c$. 
			\itemch $a=3c$. 
			\itemch $P=a+b+5c=12$. 
	\end{itemchoice}}
\end{ex}

% nguồn: Ok_dungsai_lop11_Lan2
\begin{ex}%[1D3V2-5]%[ma:4d2604a6]
	Cho $a$, $b$, $c$ là các số thực khác $0$. Biết  giới hạn $L=\lim\limits_{x\to-\infty }\dfrac{\sqrt{x^2-3x}+ax}{bx-1}$.
	\choiceTF 
		 {\True $L=3$ khi $\dfrac{a-1}b=3$}
		 {$L=6$ khi $\dfrac{a-1}b=4$}
		{\True  $L=2$ khi $ a+2b=1$}
		{\True  $L=1$ khi $ a+b=1$}
	\loigiai{Ta  có
		\begin{align*}
			\lim\limits_{x\to-\infty }\dfrac{\sqrt{x^2-3x}+ax}{bx-1}&=\lim\limits_{x\to-\infty }\dfrac{x^2-3x-{{\left(ax\right)}^2}}{\left(bx-1\right)\left(\sqrt{x^2-3x}-ax\right)}\\&=\lim\limits_{x\to-\infty }\dfrac{x\left[\left(1-{a^2}\right)x-3\right]}{\left(bx-1\right)\left(\sqrt{x^2-3x}-ax\right)}\\&=\lim\limits_{x\to-\infty }\dfrac{\left(1-{a^2}\right)-\dfrac3x}{\left(b-\dfrac1x\right)\left(-\sqrt{1-\dfrac3x}-a\right)}\\&=\dfrac{\left(1-{a^2}\right)}{b\left(-1-a\right)}=\dfrac{a-1}b. 
		\end{align*}
		\begin{itemchoice}
			 \itemch Ta có $L=\dfrac{a-1}b$. Khi đó $L=3$ thì $\dfrac{a-1}b=3$.
			\itemch $L=6$ khi $\dfrac{a-1}b=4$. 
			\itemch Khi $L=2$ thì $\dfrac{a-1}b=2\Leftrightarrow  a-1=2b\Leftrightarrow a-2b=1$.
			\itemch Khi $L=1$ thì $\dfrac{a-1}b=1\Leftrightarrow a-1=b\Leftrightarrow a-b=1$. 
	\end{itemchoice}}
\end{ex}

% nguồn: Ok_dungsai_lop11_Lan2
\begin{ex}%[1D3V2-5]%[ma:19ad8467]
	 Cho $L=\lim\limits_{x\to-\infty }\left(\sqrt{x^2+ax+5}+x\right)$.
	\choiceTF 
		{\True $L=5$ khi $ a=-10$} 
		{$L>0$ khi $  a>0$}
		 {\True $L<0$ khi $ a>0$}
		{\True $L=-1$ thì $a$ là một nghiệm của  phương trình $x^2-3x+2=0$}
	\loigiai{Ta có 
		\begin{align*}
			L&=\lim\limits_{x\to-\infty  }\left(\sqrt{x^2+ax+5}+x\right)\\&=\lim\limits_{x\to-\infty }\dfrac{ax+5}{\sqrt{x^2+ax+5}-x}\\&=\lim\limits_{x\to-\infty }\dfrac{x\left(a+\dfrac5x\right)}{|x |\sqrt{1+\dfrac ax+\dfrac5x^2}-x}\\&=\lim\limits_{x\to-\infty }\dfrac{a+\dfrac5x}{-\sqrt{1+\dfrac ax+\dfrac5x^2}-1}\\&=-\dfrac a2.
		\end{align*}
		\begin{itemchoice}
			 \itemch Khi đó $L=5\Leftrightarrow-\dfrac a2=5\Leftrightarrow a=-10$.
			\itemch Khi  đó $L>0\Leftrightarrow-\dfrac a2>0\Leftrightarrow a<0$.
			\itemch Khi đó $L<0\Leftrightarrow-\dfrac  a2<0\Leftrightarrow a>0$.
			\itemch Khi đó $L=-1\Leftrightarrow-\dfrac  a2=-1\Leftrightarrow a=2$.\\
			Ta  có $2^2-3\cdot2+2=0$ nên $a$ là một nghiệm của phương trình $x^2-3x+2=0$.
	\end{itemchoice}}
\end{ex}

% nguồn: Ok_dungsai_lop11_Lan2
\begin{ex}%[1D3V2-5]%[ma:717f8cd2]
	Cho $\lim\limits_{x\to-\infty }\dfrac{a\sqrt{x^2+1}+2023}{x+2024}=\dfrac12$; $\lim\limits_{x\to+\infty }\left(\sqrt{x^2+bx+1}-x\right)=2$. 
	 \choiceTF
		 {$ a>0$}
		{\True  $ b>0$}
		{$ a>b$} 
		 {\True $P=4a+b=2$}
	\loigiai{\begin{itemchoice}
		\itemch Ta  có
		\begin{align*}
			\lim\limits_{x\to-\infty }\dfrac{a\sqrt{x^2+1}+2023}{x+2024}&=\lim\limits_{x\to-\infty  }\dfrac{x\left(-a\sqrt{1+\dfrac1x^2}+\dfrac{2023}x\right)}{x\left(1+\dfrac{2024}x\right)}\\&=\lim\limits_{x\to-\infty }\dfrac{-a\sqrt{1+\dfrac1x^2}+\dfrac{2023}x}{1+\dfrac{2024}x}\\&=-a.
		\end{align*}
		Nên $-a=\dfrac12$ $\Leftrightarrow  a=-\dfrac12$.
		 \itemch Ta có
		\begin{align*}
			\lim\limits_{x\to+\infty }\left(\sqrt{x^2+bx+1}-x\right)&=\lim\limits_{x\to+\infty  }\dfrac{\left(\sqrt{x^2+bx+1}-x\right)\left(\sqrt{x^2+bx+1}+x\right)}{\sqrt{x^2+bx+1}+x}\\&=\lim\limits_{x\to+\infty }\dfrac{bx+1}{x\left(\sqrt{1+\dfrac bx+\dfrac1x^2}+1\right)}\\&=\lim\limits_{x\to+\infty }\dfrac{x\left(b+\dfrac1x\right)}{x\left(\sqrt{1+\dfrac bx+\dfrac1x^2}+1\right)}\\&=\lim\limits_{x\to+\infty }\dfrac{b+\dfrac1x}{\sqrt{1+\dfrac bx+\dfrac1x^2}+1}\\&=\dfrac b2.
		\end{align*}
		Nên  $\dfrac b2=2$ $\Leftrightarrow b=4$.
		 \itemch Có $ a=-\dfrac12$ và $ b=4$.
		\itemch Vậy $P=4\left(-\dfrac12\right)+4=2$. 
\end{itemchoice}}
\end{ex}

% nguồn: Ok_dungsai_lop11_Lan2
\begin{ex}%[1D3V2-5]%[ma:91188a11]
	Biết giới hạn $\lim\limits_{x \to 3}\dfrac{x+1-\sqrt{5 x+1}}{x-\sqrt{4 x-3}}=\dfrac{a}{b}$, với $a, b \in \mathbb{Z}$, $b>0$ và $\dfrac{a}{b}$  là phân số tối giản.
	\choiceTF 
		{\True $a-b=1$} 
		 {$a+b=15$}
		{$a-2 b=-9$} 
		{\True $2 a+b=26$} 
	 \loigiai{
		Ta có 
		\begin{align*}
			\lim\limits_{x \to 3}\dfrac{x+1-\sqrt{5  x+1}}{x-\sqrt{4 x-3}}
			&=\lim\limits_{x \to 3}\left[\dfrac{x+\sqrt{4 x-3}}{x+1+\sqrt{5 x+1}} \cdot \dfrac{(x+1)^2-(5 x+1)}{x^2-4 x+3}\right] \\ 
			&=\lim\limits_{x \to 3}\left[\dfrac{x+\sqrt{4  x-3}}{x+1+\sqrt{5 x+1}} \cdot \dfrac{x^2-3 x}{x^2-4 x+3}\right] \\
			&=\lim\limits_{x \to 3}\left[\dfrac{x+\sqrt{4  x-3}}{x+1+\sqrt{5 x+1}} \cdot \dfrac{x}{x-1}\right] \\
			&=\dfrac{6}{8} \cdot \dfrac{3}{2}  = \dfrac{9}{8}.
		\end{align*}
		Suy ra $a=9$, $b=8$.  Do đó $a-b=1$.
		\begin{itemchoice}
			\itemch \textbf{Đúng.} Ta  có $a=9, b=8 \Rightarrow a-b=1$.
			\itemch  \textbf{Sai.} Ta có $a=9, b=8 \Rightarrow a+b=17$.
			\itemch \textbf{Sai.} Ta có $a=9, b=8 \Rightarrow  a-2 b=-7$.
			\itemch  \textbf{Đúng.} Ta có $a=9, b=8 \Rightarrow 2 a+b=26$.
		\end{itemchoice}
	 }
\end{ex}

% nguồn: Ok_dungsai_lop11_Lan2
\begin{ex}%[1D3V2-5]%[ma:55abd14c]
	 Cho $a$, $b$ là hai số nguyên thỏa $2 a-5 b=-8$ và $\lim\limits_{x \to 0}\dfrac{\sqrt[3]{a x+1}-\sqrt{1-b x}}{x}=4$.
	\choiceTF 
		 {$|a|\leq 5$}
		{\True $a-b>1$} 
		{\True $a^2+b^2>50$} 
		{\True  $a+b>9$}
	 \loigiai{
		Ta có 
		\begin{align*}
			\lim\limits_{x \to 0}\dfrac{\sqrt[3]{a x+1}-\sqrt{1-b x}}{x} 
			&=\lim\limits_{x \to  0}\dfrac{\sqrt[3]{a x+1}-1+1-\sqrt{1-b x}}{x} \\
			&=\lim\limits_{x  \to 0}\left(\dfrac{\sqrt[3]{a x+1}-1}{x}+\dfrac{1-\sqrt{1-b x}}{x}\right) \\
			&=\lim\limits_{x \to 0}\left(\dfrac{a x+1-1}{x\left[(\sqrt[3]{a x+1})^2+\sqrt[3]{a x+1}+1\right]}+\dfrac{1-(1-b x)}{x(1+\sqrt{1-b x})}\right) \\ 
			&=\lim\limits_{x  \to 0}\left(\dfrac{a}{(\sqrt[3]{a x+1})^2+\sqrt[3]{a x+1}+1}+\dfrac{b}{1+\sqrt{1-b x}}\right) \\
			 &=\dfrac{a}{3}+\dfrac{b}{2}.
		\end{align*}
		Theo giả thiết 
		 \[
		\lim_{x  \to 0}\dfrac{\sqrt[3]{a x+1}-\sqrt{1-b x}}{x}=4
		 \Rightarrow
		 \dfrac{a}{3}+\dfrac{b}{2}=4
		 \Leftrightarrow
		 2 a+3 b=24.
		 \]
		Suy ra hệ 
		 \[
		\heva{&2 a-5 b=-8 \\ &2 a+3  b=24}
		 \Leftrightarrow
		 \heva{&a=6 \\ &b=4.}
		 \]
		\begin{itemchoice}
			\itemch \textbf{Sai.} Vì  $|a|=6>5$.
			\itemch \textbf{Đúng.} Vì $a-b=6-4=2>1$. 
			\itemch \textbf{Đúng.} Vì $a^2+b^2=6^2+4^2=52>50$. 
			 \itemch \textbf{Đúng.} Vì $a+b=6+4=10>9$.
		\end{itemchoice}
	 }
\end{ex}


%%%%%%%%%%%%%%%% TRẢ LỜI NGẮN %%%%%%%%%%%%%%%%

% nguồn: Ok_TraLoiNgan_Lop11_Lan1
\begin{ex}%[1D3H2-5]%[ma:3e8f4427]
	Tính giới hạn $\lim\limits_{x\to 2}\dfrac{\sqrt{5x-1}-\sqrt{9x-2}+1}{x-2}$  kết quả làm tròn đến hàng phần mười.
	\shortans{-0,3} 
	 \loigiai{
		\begin{eqnarray*}
			 \lim\limits_{x\to 2}\dfrac{\sqrt{5x-1}-\sqrt{9x-2}+1}{x-2}
			&=&\lim\limits_{x\to  2}\dfrac{\sqrt{5x-1}-3+4-\sqrt{9x-2}}{x-2}\\ &=&\lim\limits_{x\to 2}\left(\dfrac{\sqrt{5x-1}-3}{x-2}+\dfrac{4-\sqrt{9x-2}}{x-2}\right)\\
			&=&\lim\limits_{x\to  2}\left[\dfrac{5x-1-9}{(x-2)(\sqrt{5x-1}+3)}+\dfrac{16-9x+2}{(x-2)(4+\sqrt{9x-2})}\right]\\
			&=&\lim\limits_{x\to  2}\left[\dfrac{5(x-2)}{(x-2)(\sqrt{5x-1}+3)}+\dfrac{9(2-x)}{(x-2)(4+\sqrt{9x-2})}\right]\\
			&=&\lim\limits_{x\to  2}\left(\dfrac{5}{\sqrt{5x-1}+3}-\dfrac{9}{4+\sqrt{9x-2}}\right)=-\dfrac{7}{24} \approx -0{,}3.
		\end{eqnarray*}
	 }
\end{ex}

% nguồn: Ok_TraLoiNgan_Lop11_Lan1
\begin{ex}%[1D3V2-5]%[ma:92b81f7e]
	Tìm giới hạn hàm số $\lim\limits_{x\to 2}\dfrac{2-\sqrt{x+2}}{x^2-3x+2}$  kết quả dạng phân số $-\dfrac{a}{b}$. Hỏi tổng $2a+3b$ bằng
	 \shortans{14}
	\loigiai{ 
		Ta  có $\lim\limits_{x\to 2}\dfrac{2-\sqrt{x+2}}{x^2-3x+2}=\dfrac{0}{0}$, để khử dạng vô định ta nhân tử và mẫu cho $(2+\sqrt{x+2})$, ta được\\
		\begin{eqnarray*}
			 \lim\limits_{x\to 2}\dfrac{2-\sqrt{x+2}}{x^2-3x+2}&=&\lim\limits_{x\to 2}\dfrac{(2-\sqrt{x+2})(2+\sqrt{x+2})}{(x-1)(x-2)(2+\sqrt{x+2})}\\
			&=&\lim\limits_{x\to  2}\dfrac{-1}{(x-1)(2+\sqrt{x+2})}=-\dfrac{1}{4}.\\
		\end{eqnarray*}
		Vậy tổng $2a+3b=2 \cdot 1 +3  \cdot 4=14$.
	 }
\end{ex}

% nguồn: Ok_TraLoiNgan_Lop11_Lan1
\begin{ex}%[1D3V2-5]%[ma:2258c176]
	Tìm giới  hạn $\lim\limits_{x\to+\infty}\left(\sqrt{x^2-x}-\sqrt[3]{x^3+1}\right)$.
	\shortans{-0,5} 
	 \loigiai{
		\begin{eqnarray*}
			\lim\limits_{x\to+\infty}\left(\sqrt{x^2-x}-\sqrt[3]{x^3+1}\right)&=& \lim\limits_{x\to+\infty}\left(\sqrt{x^2-x}-x\right)+\lim\limits_{x\to+\infty}\left(x-\sqrt[3]{x^3+1}\right)\\ 
			 &=&\lim\limits_{x\to+\infty}\dfrac{x^2-x-x^2}{\sqrt{x^2-x}+x}+\lim\limits_{x\to+\infty}\dfrac{x^3-x^3-1}{x^2+x\sqrt[3]{x^3+1}+\sqrt[3]{\left(x^3+1\right)^2}}\\
			&=&\lim\limits_{x\to+\infty}\dfrac{-x}{\sqrt{x^2-x}+x}+\lim\limits_{x\to+\infty}\dfrac{-1}{x^2+x\sqrt[3]{x^3+1}+\sqrt[3]{\left(x^3+1\right)^2}} 
		\end{eqnarray*}
		Xét $A=\lim\limits_{x\to+\infty}\dfrac{-x}{\sqrt{x^2-x}+x}= \lim\limits_{x\to+\infty}\dfrac{-x}{\sqrt{x^2\left(1-\dfrac{1}{x}\right)}+x}\\ 
		=\lim\limits_{x\to+\infty}\dfrac{-x}{x\sqrt{1-\dfrac{1}{x}}+x}=\lim\limits_{x\to+\infty}\dfrac{-1}{\sqrt  1+1}=-\dfrac{1}{2}$.\\
		Xét $B=\lim\limits_{x\to+\infty}\dfrac{-1}{x^2+x\sqrt[3]{x^3+1}+\sqrt[3]{\left(x^3+1\right)^2}}=0$,\\ 
		do  $\lim\limits_{x\to+\infty}\left(x^2+x\sqrt[3]{x^3+1}+\sqrt[3]{\left(x^3+1\right)^2}\right)=+\infty $\\
		Suy ra $\lim\limits_{x\to+\infty}\left(\sqrt{x^2-x}-\sqrt[3]{x^3+1}\right)=A+B=-\dfrac{1}{2}=-0{,}5$. 
	 }
\end{ex}

% nguồn: Ok_TraLoiNgan_Lop11_Lan1
\begin{ex}%[1D3V2-5]%[ma:35102f79]
	Tính giới hạn $\lim\limits_{x \to 2} \dfrac{\sqrt{x+7}-3}{x-2}$ (Kết quả làm tròn đến hàng thập  phân thứ $2$).
	\shortans{0,17} 
	 \loigiai{
		Ta có $\lim\limits_{x \to 2} \dfrac{\sqrt{x+7}-3}{x-2}=\lim\limits_{x  \to 2}\dfrac{x+7-9}{(x-2)\left(\sqrt{x+7}+3\right)}=\lim\limits_{x \to 2}=\dfrac{1}{\sqrt{x+7}+3}=\dfrac{1}{6}\approx 0{,}17$.
	 }
\end{ex}

% nguồn: Ok_TraLoiNgan_Lop11_Lan1
\begin{ex}%[1D3V2-5]%[ma:6a535f00]
	Tính giới hạn của hàm số sau  $\lim\limits_{x\to-2}\dfrac{\sqrt[3]{2x+12}+x}{x^2+2x}$ kết quả dạng phân số $-\dfrac{a}{b}$. Hỏi tổng $a+2b$ bằng
	 \shortans{17}
	 \loigiai{
		\begin{eqnarray*}
			 \lim\limits_{x\to-2}\dfrac{\sqrt[3]{2x+12}+x}{x^2+2x}&=&\lim\limits_{x\to-2}\dfrac{(\sqrt[3]{2x+12}+x)\left(\sqrt[3]{(2x+12)^2}-x\sqrt[3]{2x+12}+x^2\right)}{x(x+2)\left(\sqrt[3]{(2x+12)^2}-x\sqrt[3]{2x+12}+x^2\right)}\\
			 &=&\lim\limits_{x\to-2}\dfrac{(x+2)\left(x^2-2x+12\right)}{x(x+2)\left(\sqrt[3]{(2x+12)^2}-x\sqrt[3]{2x+12}+x^2\right)}\\
			&=&\lim\limits_{x\to-2}\dfrac{x^2-2x+12}{\left(\sqrt[3]{(2x+12)^2}-x\sqrt[3]{2x+12}+x^2\right)x}=-\dfrac{5}{6}. 
		\end{eqnarray*}
		Vậy tổng $a+2b=5+2\cdot  6=17$.
	} 
\end{ex}

% nguồn: Ok_TraLoiNgan_Lop11_Lan1
\begin{ex}%[1D3V2-5]%[ma:51fc0a70]
	Tính giới  hạn $\lim\limits_{x\to 1}\dfrac{\sqrt{3x+1}-2}{2x^2-3x+1}$
	 \shortans{0,75}
	 \loigiai{
		\begin{eqnarray*}
			\lim\limits_{x\to 1}\dfrac{\sqrt{3x+1}-2}{2x^2-3x+1} 
			&=&\lim\limits_{x\to 1}\dfrac{3x+1-4}{\left(2x^2-3x+1\right)(\sqrt{3x+1}+2)}\\ 
			 &=&\lim\limits_{x\to 1}\dfrac{3(x-1)}{(x-1)(2x-1)(\sqrt{3x+1}+2)}\\
			&=&\lim\limits_{x\to 1}\dfrac{3}{(2x-1)(\sqrt{3x+1}+2)}=\dfrac{3}{4}  = 0{,}75.
		\end{eqnarray*}
	} 
\end{ex}

% nguồn: Ok_TraLoiNgan_Lop11_Lan1
\begin{ex}%[1D3V2-5]%[ma:a2dde504]
	Tính giới hạn $\lim\limits_{x\to-1}\dfrac{3x+\sqrt{4-5x}}{-2x^2+3x+5}$ kết quả làm tròn đến hàng  phần trăm.
	 \shortans{0,31}
	\loigiai{ 
		\begin{eqnarray*}
			 \lim\limits_{x\to-1}\dfrac{3x+\sqrt{4-5x}}{-2x^2+3x+5}
			 &=&\lim\limits_{x\to-1}\dfrac{9x^2-4+5x}{\left(-2x^2+3x+5\right)(3x-\sqrt{4-5x})}\\
			&=&\lim\limits_{x\to-1}\dfrac{(x+1)(9x-4)}{(x+1)(-2x+5)(3x-\sqrt{4-5x})}\\ 
			 &=&\lim\limits_{x\to-1}\dfrac{9x-4}{(-2x+5)(3x-\sqrt{4-5x})}=\dfrac{13}{42} \approx 0{,}31.
		\end{eqnarray*}
	} 
\end{ex}

% nguồn: Ok_TraLoiNgan_Lop11_Lan1
\begin{ex}%[1D3V2-5]%[ma:130a9921]
	Tìm giới hạn $\lim\limits_{x\to  0}\dfrac{\sqrt[3]{1+x^2}-1}{x^2}$ (kết quả làm tròn đến hàng phần trăm).
	\shortans{0,33} 
	\loigiai{ 
		\begin{eqnarray*}
			 \lim\limits_{x\to 0}\dfrac{\sqrt[3]{1+x^2}-1}{x^2}&=&\lim\limits_{x\to 0}\dfrac{1+x^2-1^3}{\left(\sqrt[3]{\left(1+x^2\right)^2}+\sqrt[3]{1+x^2}+1\right){x^2}}\\
			&=&\lim\limits_{x\to 0}\dfrac{1}{\sqrt[3]{(1+x^2)^2}+\sqrt[3]{1+x^2}+1}=\dfrac{1}{3} \approx 0{,}33. 
		\end{eqnarray*}
	} 
\end{ex}

% nguồn: Ok_TraLoiNgan_Lop11_Lan1
\begin{ex}%[1D3V2-5]%[ma:16d3201e]
	Tìm giới hạn $\lim\limits_{x\to 1}\dfrac{\sqrt[3]{x+7}-\sqrt{x+3}}{x^2-3x+2}$ (kết quả  làm tròn đến hàng phần trăm).
	\shortans{0,17} 
	\loigiai{ 
		 $\lim\limits_{x\to 1}\dfrac{\sqrt[3]{x+7}-\sqrt{x+3}}{x^2-3x+2}= \lim\limits_{x\to 1}\dfrac{\sqrt[3]{x+7}-2}{(x-1)(x-2)}+\lim\limits_{x\to 1}\dfrac{2-\sqrt{x+3}}{(x-1)(x-2)}$.\\
		Ta  có
		\begin{eqnarray*}
			 \lim\limits_{x\to 1}\dfrac{\sqrt[3]{x+7}-2}{(x-1)(x-2)}&=& \lim\limits_{x\to 1}\dfrac{x+7-2^3}{(x-1)(x-2)\left[\sqrt[3]{(x+7)^2}+2\sqrt[3]{x+7}+4\right]}\\
			&=&\lim\limits_{x\to 1}\dfrac{1}{(x-2)\left(\sqrt[3]{(x+7)^2}+2\sqrt[3]{x+7}+4\right)}= -\dfrac{1}{12}.\\ 
			\lim\limits_{x\to 1}\dfrac{2-\sqrt{x+3}}{(x-1)(x-2)}&=& \lim\limits_{x\to 1}\dfrac{2^2-(x+3)}{(x-1)(x-2)(2+\sqrt{x+3})}\\ 
			&=& \lim\limits_{x\to 1}\dfrac{-1}{(x-2)(2+\sqrt{x+3})}=\dfrac{1}{4}. 
		\end{eqnarray*}
		Vậy $\lim\limits_{x\to 1}\dfrac{\sqrt[3]{x+7}-\sqrt{x+3}}{x^2-3x+2}= -\dfrac{1}{12}+\dfrac{1}{4}=\dfrac{1}{6}  \approx 0{,}17$.
	 }
\end{ex}

% nguồn: Ok_TraLoiNgan_Lop11_Lan1
\begin{ex}%[1D3V2-5]%[ma:42db42fe]
	Tính giới hạn $\lim\limits_{x\to 1}\dfrac{\sqrt{x+3}+\sqrt{2x+7}-5}{2x-2}$ kết quả làm tròn đến hàng phần  trăm.
	 \shortans{0,29}
	 \loigiai{
		\begin{eqnarray*}
			\lim\limits_{x\to 1}\dfrac{\sqrt{x+3}+\sqrt{2x+7}-5}{2x-2} 
			 &=&\lim\limits_{x\to 1}\dfrac{\sqrt{x+3}-2+\sqrt{2x+7}-3}{2x-2}\\
			 &=&\lim\limits_{x\to 1}\left[\dfrac{x+3-4}{(2x-2)(\sqrt{x+3}+2)}+\dfrac{2x+7-9}{(2x-2)(\sqrt{2x+7}+3)}\right]\\
			&=&\lim\limits_{x\to 1}\left[\dfrac{x-1}{2(x-1)(\sqrt{x+3}+2)}+\dfrac{2x-2}{(2x-2)(\sqrt{2x+7}+3)}\right]\\ 
			&=&\lim\limits_{x\to 1}\left[\dfrac{1}{2(\sqrt{x+3}+2)}+\dfrac{1}{\sqrt{2x+7}+3}\right] 
			 =\dfrac{7}{24} \approx 0{,}29.
		\end{eqnarray*}
	 }
\end{ex}

% nguồn: Ok_TraLoiNgan_Lop11_Lan1
\begin{ex}%[1D3V2-5]%[ma:200a36d3]
	Giới hạn của hàm số sau $\displaystyle\lim\limits_{x\to 1}\dfrac{\sqrt{2x^2+2}-\sqrt{x^2+3x}}{x-1}$ có kết quả là $-\dfrac{a}{b}$, với $\dfrac{a}{b}$  là phân số tối giản. Tính $a-b$.
	 \shortans{-3}
	 \loigiai{
	Ta  có
	\begin{eqnarray*}
		 & & \displaystyle\lim\limits_{x\to 1}\dfrac{\sqrt{2x^2+2}-\sqrt{x^2+3x}}{x-1}\\
		 &=&\displaystyle\lim\limits_{x\to 1}\dfrac{x^2-3x+2}{(x-1)(\sqrt{2x^2+2}+\sqrt{x^2+3x})}\\
		&=&\displaystyle\lim\limits_{x\to 1}\dfrac{x-2}{\sqrt{2x^2+2}+\sqrt{x^2+3x}}=-\dfrac{1}{4}. 
	\end{eqnarray*}
	Suy ra: $a=1,b=4$. 
	Vậy $a-b=-3$. 
	 }
\end{ex}

% nguồn: Ok_TraLoiNgan_Lop11_Lan1
\begin{ex}%[1D3V2-5]%[ma:dbc354b4]
	Cho $L=\displaystyle\lim_{x\to 1}\dfrac{\sqrt{x+3}-3x+1}{x^2+x-2}=\dfrac{-a}{b}$  với $a$, $b$ là các số nguyên dương và $\dfrac{a}{b}$ là phân số tối giản. Tính $a+b$.
	\shortans{23} 
	 \loigiai{
		$\displaystyle\lim_{x\to 1}\dfrac{\sqrt{x+3}-3x+1}{x^2+x-2}=\displaystyle\lim_{x\to 1}\dfrac{\sqrt{x+3}-2-3(x-1)}{(x-1)(x+2)}=\displaystyle\lim_{x\to 1}\left[\dfrac{\sqrt{x+3}-2}{(x-1)(x+2)}-\dfrac{3}{x+2}\right]$\\ 
		$=\displaystyle\lim_{x\to 1}\left[\dfrac{(x-1)}{(x-1)(x+2)(\sqrt{x+3}+2)}-\dfrac{3}{x+2}\right]=\displaystyle\lim_{x\to 1}\left[\dfrac{1}{(x+2)(\sqrt{x+3}+2)}-\dfrac{3}{x+2}\right]=-\dfrac{11}{12}$.\\ 
		Vậy $a=11$,  $b=12$ suy ra $a+b=23$.
	} 
\end{ex}

% nguồn: Ok_TraLoiNgan_Lop11_Lan2
\begin{ex}%[1D3V2-5]%[ma:8ee08da6]
	Biết $\lim\limits_{x\to+\infty }\left(\sqrt{3x^2+2023x+2}-n\sqrt3\right)=\dfrac{a\sqrt b}c$, ở đó $ a,c$ là các số nguyên dương nguyên tố cùng nhau, $b$ là số nguyên tố. Khi đó hãy tính giá  trị của $ a+b+c$.
	\shortans{2032} 
	\loigiai{Ta có $\dfrac{a\sqrt b}c=\lim\limits_{x\to+\infty  }\dfrac{2023x+2}{\sqrt{3x^2+2023x+2}+x\sqrt3}=\lim\limits_{x\to+\infty }\dfrac{2023+\dfrac2x}{\sqrt{3+\dfrac{2023}x+\dfrac2x^2}+\sqrt3}=\dfrac{2023\sqrt3}6.$\\
	Vậy $  a+b+c=2023+3+6=2032$.}
\end{ex}

% nguồn: Ok_TraLoiNgan_Lop11_Lan2
\begin{ex}%[1D3V2-5]%[ma:3f9e2287]
	Cho $\lim\limits_{x  \to 3}\dfrac{x+1-\sqrt{5 x+1}}{x-\sqrt{4 x-3}}=\dfrac{a}{b}$, với $a, b \in \mathbb{Z}$, $b>0$ và $\dfrac{a}{b}$ là phân số tối giản.
	Giá trị của $a-b$  bằng
	\shortans{1} 
	 \loigiai{
		Ta  có
		\begin{align*}
			\lim\limits_{x \rightarrow 3}  \dfrac{x+1-\sqrt{5 x+1}}{x-\sqrt{4 x-3}}
			&=\lim\limits_{x  \rightarrow 3} \dfrac{\left[(x+1)^2-(5 x+1)\right](x+\sqrt{4 x-3})}{\left[x^2-(4 x-3)\right](x+1+\sqrt{5 x+1})} \\
			&=\lim\limits_{x  \rightarrow 3} \dfrac{\left(x^2-3 x\right)(x+\sqrt{4 x-3})}{\left(x^2-4 x+3\right)(x+1+\sqrt{5 x+1})} \\
			&=\lim\limits_{x \rightarrow 3} \dfrac{x(x-3)(x+\sqrt{4 x-3})}{(x-1)(x-3)(x+1+\sqrt{5 x+1})}  \\
			 &=\lim\limits_{x \rightarrow 3} \dfrac{x(x+\sqrt{4 x-3})}{(x-1)(x+1+\sqrt{5 x+1})} \\
			&=\dfrac{3 \cdot 6}{2 \cdot 8}=\dfrac{9}{8}. 
		\end{align*}
		Khi đó $a=9$, $b=8$. Do vậy $a-b=1$. 
	 }
\end{ex}

% nguồn: Ok_TraLoiNgan_Lop11_Lan2
\begin{ex}%[1D3V2-5]%[ma:c50bc383]
	 Cho $\lim\limits_{x\to+\infty }\left(\dfrac{x^2+3x+1}{x+1} \texttt{+}ax+b\right)=1$. Tính giá trị của biểu thức $T=2024a-4049b$.
	\shortans{2025} 
	\loigiai{Ta  có
		\begin{align*}
			\lim\limits_{x\to+\infty }\left(\dfrac{x^2+3x+1}{x+1} \texttt{+}ax+b\right)=1&\Leftrightarrow\lim\limits_{x\to+\infty }\left(\dfrac{\left(a+1\right)x^2+\left(a+b+3\right)x+b+1}{x+1}\right)=1\\&\Leftrightarrow\lim\limits_{x\to+\infty  }\left(\dfrac{\left(a+1\right)x+\left(a+b+3\right)+\dfrac{b+1}x}{1+\dfrac1x}\right)=1\\&\Leftrightarrow\heva{&a+1=0\\&a+b+3=1\\
			 }\Leftrightarrow\heva{&a=-1\\&b=-1\\
			}\\&\Rightarrow  T=2024a-4049b=2024\cdot(-1)-4049(-1)=2025.
	\end{align*}}
\end{ex}

% nguồn: Ok_TraLoiNgan_Lop11_Lan2
\begin{ex}%[1D3V2-5]%[ma:e431326e]
	Biết rằng  giới hạn
	$$\lim\limits_{x\to 2}\dfrac{\sqrt{5x-1}-\sqrt{9x-2}+1}{x-2}=\lim\limits_{x\to 2}\left(\dfrac a{b\cdot\sqrt{5x-1}+3}-\dfrac c{4+d\cdot\sqrt{9x-2}}\right)$$ 
	 với $a$; $b$; $c$; $d$ là các số nguyên dương. Tính tổng các số $a$; $b$; $c$; $d$.
	\shortans{16} 
	\loigiai{Ta  có
		\begin{align*}
			 \lim\limits_{x\to 2}\dfrac{\sqrt{5x-1}-\sqrt{9x-2}+1}{x-2}&=\lim\limits_{x\to 2}\dfrac{\sqrt{5x-1}-3+4-\sqrt{9x-2}}{x-2}\\&=\lim\limits_{x\to 2}\left(\dfrac{\sqrt{5x-1}-3}{x-2}+\dfrac{4-\sqrt{9x-2}}{x-2}\right)\\&=\lim\limits_{x\to 2}\left[\dfrac{5x-1-9}{\left(x-2\right)\left(\sqrt{5x-1}+3\right)}+\dfrac{16-9x+2}{\left(x-2\right)\left(4+\sqrt{9x-2}\right)}\right]\\&=\lim\limits_{x\to 2}\left[\dfrac{5\left(x-2\right)}{\left(x-2\right)\left(\sqrt{5x-1}+3\right)}+\dfrac{9\left(2-x\right)}{\left(x-2\right)\left(4+\sqrt{9x-2}\right)}\right]\\&=\lim\limits_{x\to 2}\left(\dfrac5{\sqrt{5x-1}+3}-\dfrac9{4+\sqrt{9x-2}}\right)\\&\Rightarrow\heva{&a=5\\&b=1\\&c=9\\&d=1}\Rightarrow S=16.
	\end{align*}}
\end{ex}

% nguồn: Ok_TraLoiNgan_Lop11_Lan2
\begin{ex}%[1D3V2-5]%[ma:c8c48e7d]
	Biết rằng giới hạn $\lim\limits_{x\to 0}\dfrac{\sqrt{x+9}+\sqrt{x+16}-7}x$ $=\lim\limits_{x\to 0}\left[\dfrac a{\sqrt{x+9}+b}+\dfrac c{\sqrt{x+16}+d}\right]$ với $a$; $b$; $c$; $d$ là các số nguyên dương. Tính tổng các số  $a$; $b$; $c$; $d$.
	 \shortans{9}
	\loigiai{Ta có 
		\begin{align*}
			\lim\limits_{x\to  0}\dfrac{\sqrt{x+9}+\sqrt{x+16}-7}x&
			=\lim\limits_{x\to 0}\left(\dfrac{\sqrt{x+9}-3}x+\dfrac{\sqrt{x+16}-4}x\right)\\&=\lim\limits_{x\to 0}\left[\dfrac{x+9-9}{x\left(\sqrt{x+9}+3\right)}+\dfrac{x+16-16}{x\left(\sqrt{x+16}+4\right)}\right]\\&=\lim\limits_{x\to 0}\left[\dfrac1{\sqrt{x+9}+3}+\dfrac1{\sqrt{x+16}+4}\right]. 
		\end{align*}
	Vậy tổng  các số $a$; $b$; $c$; $d$ là $9$.}
\end{ex}

% nguồn: Ok_TraLoiNgan_Lop11_Lan2
\begin{ex}%[1D3V2-5]%[ma:4dcc4ba2]
	Giá trị  của $\lim\limits_{x\to 1}\dfrac{\sqrt{3x+1}-2}{2x^2-3x+1}$ bằng bao nhiêu?
	\shortans{0,75} 
	\loigiai{Ta  có
		\begin{align*}
			\lim\limits_{x\to 1}\dfrac{\sqrt{3x+1}-2}{2x^2-3x+1}&=\lim\limits_{x\to 1}\dfrac{3x+1-4}{\left(2x^2-3x+1\right)\left(\sqrt{3x+1}+2\right)}\\ 
			 &=\lim\limits_{x\to 1}\dfrac{3\left(x-1\right)}{\left(x-1\right)\left(2x-1\right)\left(\sqrt{3x+1}+2\right)}\\
			 &=\lim\limits_{x\to 1}\dfrac3{\left(2x-1\right)\left(\sqrt{3x+1}+2\right)}\\
			 &=\dfrac34=0{,}75.
	\end{align*}}
\end{ex}

% nguồn: Ok_TraLoiNgan_Lop11_Lan1
\begin{ex}%[1D3C2-5]%[ma:bc29c2eb]
	Cho hàm số $y=f(x)$ xác định trên $\mathbb {R}$ thỏa  mãn $\lim\limits_{x\to 2} \dfrac{f(x)-16}{x-2}=12$ khi đó giới hạn \break $\lim\limits_{x\to 2} \dfrac{\sqrt{2f(x)-16}-4}{x^2+x-6}$ bằng bao nhiêu \textit{(kết quả là số thập phân làm tròn hàng phần mười)}?\shortans{0,6}
	\loigiai{  Ta có $\lim\limits_{x\to 2} \dfrac{f(x)-16}{x-2}=12\Rightarrow \lim\limits_{x\to 2} f(x)= \lim\limits_{x\to 2} [12(x-2)+16]= \lim\limits_{x\to 2} (12x-8)$.\\
		Từ đó ta có \begin{eqnarray*}
		 \lim\limits_{x\to 2} \dfrac{\sqrt{2f(x)-16}-4}{x^2+x-6}
		&=&\lim\limits_{x\to  2} \dfrac{\sqrt{2(12x-8)-16}-4}{x^2+x-6}\\
		&=&\lim\limits_{x\to  2} \dfrac{\sqrt{24x-32}-4}{x^2+x-6}\\
		&=&\lim\limits_{x\to  2} \dfrac{24(x-2)}{(x-2)(x+3)\left(\sqrt{24x-32}+4\right)}\\
		&=&\lim\limits_{x\to 2}  \dfrac{24}{(x+3)\left(\sqrt{24x-32}+4\right)}\\
		&=&\dfrac  {3}{5}=0{,}6.
	\end{eqnarray*}
} 
\end{ex}


%%%%%%%%%%%%%%%% TỰ LUẬN %%%%%%%%%%%%%%%%

% nguồn: Ok_TuanLuan_Toan11_Lan2
\begin{ex}%[1D3H2-5]%[ma:0203be15]
	$\lim\limits_{x\to+\infty}x\left(\sqrt{x^2+5x+4}-\sqrt{x^2+5x-2}\right)$ bằng 
	 \loigiai{
		\allowdisplaybreaks 
		\begin{eqnarray*}
			\lim\limits_{x\to+\infty}x\left(\sqrt{x^2+5x+4}-\sqrt{x^2+5x-2} \right)&=&\lim\limits_{x\to+\infty}  \dfrac{6x}{\sqrt{x^2+5x+4}+\sqrt{x^2+5x-2}}\\
			&=&\lim\limits_{x\to+\infty}  \dfrac{6x}{x\left(\sqrt{1+\dfrac{5}{x}+\dfrac{4}{x^2}}+\sqrt{1+\dfrac{5}{x}-\dfrac{2}{x^2}}\right)}\\
			 &=&3.
		\end{eqnarray*}
	 }
\end{ex}

% nguồn: Ok_TuanLuan_Toan11_Lan2
\begin{ex}%[1D3H2-5]%[ma:3d6cc5af]
	Tính giới hạn  $\lim\limits_{x\to 2}\dfrac{\sqrt{x+2}-2}{x-2}$.
	\loigiai{ 
	 $ \lim\limits_{x\to 2}\dfrac{\sqrt{x+2}-2}{x-2}=\lim\limits_{x\to 2}\dfrac{x-2}{\left(x-2\right)\left(\sqrt{x+2}+2\right)}= \lim\limits_{x\to 2}\dfrac{1}{\sqrt{x+2}+2}=\dfrac{1}{4}$.
	} 
\end{ex}

% nguồn: Ok_TuanLuan_Toan11_Lan2
\begin{ex}%[1D3H2-5]%[ma:d32cac54]
	Biết $ \lim\limits_{x\to  0}\dfrac{\sqrt{3x+1}-1}{x}=\dfrac{a}{b}$, trong đó $a$, $b$ là các số nguyên dương và phân số $\dfrac{a}{b}$ tối giản. Tính giá trị biểu thức $P=a^2+b^2$.
	 \loigiai{
	Ta có: $ \lim\limits_{x\to 0}\dfrac{\sqrt{3x+1}-1}{x}= \lim\limits_{x\to 0}\dfrac{3x+1-1}{x\left( \sqrt{3x+1}+1  \right)}= \lim\limits_{x\to 0}\dfrac{3}{\sqrt{3x+1}+1}=\dfrac{3}{2}$.\\
	Do  đó, $a=3$, $b=2$.\\
	 Vậy $P={a^2}+{{b}^{2}}=13$.
	} 
\end{ex}

% nguồn: Ok_TuanLuan_Toan11_Lan2
\begin{ex}%[1D3H2-5]%[ma:f32f9294]
	 Tính các giới hạn sau đây
	\begin{multicols}{2}
		\begin{enumerate}
			\item $\displaystyle\lim_{x\to 0} \dfrac{\sqrt{1+2x}-1}{2x}$; 
			\item $\displaystyle\lim_{x\to 0} \dfrac{4x}{\sqrt{9+x}-3}$; 
			\item $\displaystyle\lim_{x\to 1}  \dfrac{\sqrt{2x+7}-3}{2-\sqrt{x+3}}$;
			\item $\displaystyle\lim_{x\to 0} \dfrac{\sqrt{1+x^2}-1}{\sqrt{x^2+16}-4}$; 
			\item  $\displaystyle\lim_{x\to 2} \dfrac{\sqrt{x+2}-\sqrt{2x}}{\sqrt{x-1}-\sqrt{3-x}}$;
			\item  $\displaystyle\lim_{x\to 1} \dfrac{\sqrt{2x+7}+x+4}{x^3-4x^2+3}$;
			\item $\displaystyle\lim_{x\to -1} \dfrac{x^2-1}{2x+\sqrt{3x^2+1}}$; 
			\item $\displaystyle\lim_{x\to 5}  \dfrac{2x-5\sqrt{x-1}}{3-\sqrt{x+4}}$;
			 \item $\displaystyle\lim_{x\to 0} \dfrac{1-\sqrt[3]{12x+1}}{4x}$.
		\end{enumerate}
	\end{multicols}
	\loigiai{ 
		\begin{enumerate}
			\item Ta  có
			\begin{eqnarray*}
				\displaystyle\lim_{x \to  0} \dfrac{\sqrt{1 + 2x}-1}{2x}
				&=& \displaystyle\lim_{x \to 0} \dfrac{2x}{2x\left(\sqrt{1+2x}+  1\right)} \\
				&=& \displaystyle\lim_{x \to 0} \dfrac{1}{\sqrt{1+2x}+1}\\ 
				&=& \dfrac{1}{2}. 
			\end{eqnarray*}
			\item Ta  có
			\begin{eqnarray*}
				\displaystyle\lim_{x  \to 0} \dfrac{4x}{\sqrt{9+x}-3}
				&=& \displaystyle\lim_{x \to  0} \dfrac{4x\left(\sqrt{9+x}+ 3\right)}{x} \\
				 &=& \displaystyle\lim_{x \to 0} 4\left(\sqrt{9+x}+3\right)\\
				 &=& 24.
			\end{eqnarray*}
			 \item Ta có
			\begin{eqnarray*}
				\displaystyle\lim_{x \to 1}  \dfrac{\sqrt{2x+7}-3}{2-\sqrt{x+3}}
				&=& \displaystyle\lim_{x \to 1} \left[\dfrac{(2x+7-9)}{\left(2- \sqrt{x+3}\right)\cdot\left(\sqrt{2x+7}+3\right)} \cdot \dfrac{\left(2+\sqrt{x+3}\right)}{\left(2+\sqrt{x+3}\right)}\right]  \\
				&=& \displaystyle\lim_{x \to  1} \dfrac{2(x-1)\left(2+\sqrt{x+ 3}\right)}{(x-1)\left(\sqrt{2x+7}+3\right)} \\
				&=& \displaystyle\lim_{x \to  1} \dfrac{2\left(2+ \sqrt{x+3}\right)}{\sqrt{2x+7}+3} \\
				 &=& \dfrac{8}{5}.
			\end{eqnarray*}
			\item Ta có 
			\begin{eqnarray*}
				 \displaystyle\lim_{x \to 0} \dfrac{\sqrt{1+x^2}-1}{\sqrt{x^2+16}-4}
				&=& \displaystyle\lim_{x \to 0} \left[\dfrac{x^2}{\left(\sqrt{1+x^2} +  1\right)} \cdot \dfrac{\sqrt{x^2+16}+4}{x^2}\right]\\
				&=& \displaystyle\lim_{x  \to 0} \dfrac{\sqrt{x^2+16}+4}{\sqrt{1+x^2} + 1} \\
				&=& 4. 
			\end{eqnarray*}
			\item Ta có 
			\begin{eqnarray*}
				\displaystyle\lim_{x \to  2} \dfrac{\sqrt{x+2}-\sqrt{2x}}{\sqrt{x-1}- \sqrt{3-x}}
				&=& \displaystyle\lim_{x  \to 2} \dfrac{(x+2-2x)\cdot\left(\sqrt{x-1} +\sqrt{3-x}\right)}{(x-1-3+x)\cdot\left(\sqrt{x + 2} + \sqrt{2x}\right)} \\
				 &=& \displaystyle\lim_{x \to 2} \dfrac{-(x-2)\cdot\left(\sqrt{x-1} +\sqrt{3-x}\right)}{2(x-2)\cdot\left(\sqrt{x+2}+\sqrt{2x}\right)} \\
				 &=& \displaystyle\lim_{x \to 2} \dfrac{-\left(\sqrt{x-1} +\sqrt{3-x}\right)}{2\left(\sqrt{x+2}+ \sqrt{2x}\right)} \\
				 &=&\dfrac{-1}{4}.
			\end{eqnarray*}
			\item Ta  có
			\begin{eqnarray*}
				\displaystyle\lim_{x \to  1} \dfrac{\sqrt{2x+7}+x-4}{x^3-4x^2+3}
				&=& \displaystyle\lim_{x \to 1} \dfrac{\sqrt{2x+7}-3+x-1}{(x-1)(x^2-3x-3)}  \\
				&=& \displaystyle\lim_{x  \to 1} \dfrac{\dfrac{2x+7-9}{\sqrt{2x+7}+3}+x-1}{(x-1)(x^2-3x-3)} \\
				&=& \displaystyle\lim_{x \to 1} \dfrac{\dfrac{2(x-1)}{\sqrt{2x+7}+3}+x-1}{(x-1)(x^2-3x-3)}  \\
				&=&  \displaystyle\lim_{x \to 1} \dfrac{\dfrac{2}{\sqrt{2x+7}+3}+1}{x^2-3x-3} \\
				 &=&\dfrac{-4}{15}.
			\end{eqnarray*}
			\item Ta  có
			\begin{eqnarray*}
				\displaystyle\lim_{x \to -1} \dfrac{x^2-1}{2x+\sqrt{3x^2+ 1}} 
				 &=& \displaystyle\lim_{x \to -1} \dfrac{(x-1)(x+1)\left(2x- \sqrt{3x^2+1}\right)}{4x^2-3x^2-1} \\
				&=& \displaystyle\lim_{x \to -1}  \dfrac{(x-1)(x+1)\left(2x- \sqrt{3x^2+1}\right)}{(x-1)(x+1)} \\
				&=& \displaystyle\lim_{x \to -1}  \left(2x-\sqrt{3x^2+1}\right)\\
				 &=& -4.
			\end{eqnarray*}
			\item  Ta có
			\begin{eqnarray*}
				\displaystyle\lim_{x \to  5} \dfrac{2x-5\sqrt{x-1}}{3-\sqrt{x+4}}
				&=& \displaystyle\lim_{x\to 5} \dfrac{\left[4x^2-25(x-1)\right]\left(3+\sqrt{x+4}\right)}{\left(9-(x+4)\right)\left(2x+5\sqrt{x-1}\right)} \\ 
				&=& \displaystyle\lim_{x\to 5} \dfrac{(x-5)(4x-5)\left(3+\sqrt{x+4}\right)}{(5-x)\left(2x+5\sqrt{x-1}\right)} \\ 
				&=& \displaystyle\lim_{x\to 5} \dfrac{(4x-5)\left(3+\sqrt{x+4}\right)}{-\left(2x+5\sqrt{x-1}\right)}  \\
				&=&\dfrac{-9}{2}. 
			\end{eqnarray*}
			 \item Ta có
			\begin{eqnarray*}
				\displaystyle\lim_{x  \to 0} \dfrac{1-\sqrt[3]{12x+1}}{4x}
				&=&  \displaystyle\lim_{x \to 0} \dfrac{1-12x-1}{4x\left(1+\sqrt[3]{12x+1}+\sqrt[3]{(12x+1)^2}\right)} \\
				&=& \displaystyle\lim_{x \to 0} \dfrac{-12x}{4x\left(1+\sqrt[3]{12x +  1}+\sqrt[3]{(12x + 1)^2}\right)} \\
				&=& \displaystyle\lim_{x  \to 0} \dfrac{-3}{1+\sqrt[3]{12x + 1}+\sqrt[3]{(12x+1)^2}} \\
				 &=& -1.
			\end{eqnarray*}
		\end{enumerate}
	} 
\end{ex}

% nguồn: Ok_TuanLuan_Toan11_Lan2
\begin{ex}%[1D3V2-5]%[ma:65e4513b]
	Biết $\lim\limits_{x\to +\infty}\left[\sqrt{4x^2-3x+1}-\left(ax+b\right)\right]=0$.  Tính $a-4b$ ta được
	 \loigiai{
		Ta có 
		 \allowdisplaybreaks
		\begin{eqnarray*}
			 \lim\limits_{x\to +\infty}\left[\sqrt{4x^2-3x+1}-\left(ax+b\right)\right]=0&\Leftrightarrow&\lim\limits_{x\to +\infty}\left[\left(\sqrt{4x^2-3x+1}-ax\right)-b\right]=0\\
			&\Leftrightarrow&\lim\limits_{x\to  +\infty}\left(\dfrac{4x^2-3x+1-a^2x^2}{\sqrt{4x^2-3x+1}+ax}-b\right)=0\\
			&\Leftrightarrow&\lim\limits_{x\to  +\infty}\left[\dfrac{(4-a^2)x^2-3x+1}{\sqrt{4x^2-3x+1}+ax}-b\right]=0\\
			&\Leftrightarrow&\heva{&4-a^2=0\\&a>0\\&\dfrac{-3}{2+a}-b=0}\\ 
			&\Leftrightarrow&\heva{&a=2\\&b=-\dfrac{3}{4}.} 
		\end{eqnarray*}
		Vậy $a-4b=5$. 
	 }
\end{ex}

% nguồn: Ok_TuanLuan_Toan11_Lan2
\begin{ex}%[1D3V2-5]%[ma:ae1e8255]
	Biết $ \lim\limits_{x\to 1}\,\dfrac{\sqrt{x^2+x+2}-\sqrt[3]{7x+1}}{\sqrt{2}(x-1)}=\dfrac{a\sqrt{2}}{b}$ với $a$, $b \in \mathbb{Z}$ và $\dfrac{a}{b}$ là phân số tối giản. Tính giá trị  của $a+b$.
	 \loigiai{
	 Ta có
	\begin{align*}
		 \lim\limits_{x\to 1}\,\dfrac{\sqrt{x^2+x+2}-\sqrt[3]{7x+1}}{\sqrt{2}(x-1)}&= \lim\limits_{x\to 1}\,\dfrac{\sqrt{x^2+x+2}-2+2-\sqrt[3]{7x+1}}{\sqrt{2}(x-1)}\\
		&=  \lim\limits_{x\to 1}\,\dfrac{\sqrt{x^2+x+2}-2}{\sqrt{2}(x-1)}+ \lim\limits_{x\to 1}\,\dfrac{2-\sqrt[3]{7x+1}}{\sqrt{2}(x-1)}\\
		 &=I+J.
	\end{align*}
	Ta  có
	\begin{align*}
		 I&= \lim\limits_{x\to 1}\dfrac{\sqrt{x^2+x+2}-2}{\sqrt{2}(x-1)}\\
		&=  \lim\limits_{x\to 1}\dfrac{x^2+x+2-4}{\sqrt{2}(x-1)\left(\sqrt{x^2+x+2}+2 \right)}\\
		 &= \lim\limits_{x\to 1}\dfrac{(x-1)(x+2)}{\sqrt{2}(x-1)\left(\sqrt{x^2+x+2}+2 \right)}\\
		&= \lim\limits_{x\to 1}\,\dfrac{x+2}{\sqrt{2}\left(\sqrt{x^2+x+2}+2  \right)}\\
		 &=\dfrac{3}{4\sqrt{2}}.
	\end{align*}
	và 
	\begin{align*}
		 J&= \lim\limits_{x\to 1}\,\dfrac{2-\sqrt[3]{7x+1}}{\sqrt{2}(x-1)}\\
		&= \lim\limits_{x\to 1}\,\dfrac{8-7x-1}{\sqrt{2}(x-1)\left[4+2\sqrt[3]{7x+1}+{{\left(  \sqrt[3]{7x+1} \right)}^{2}} \right]}\\
		&= \lim\limits_{x\to 1}\,\dfrac{-7}{\sqrt{2}\left[4+2\sqrt[3]{7x+1}+{{\left(\sqrt[3]{7x+1}  \right)}^{2}} \right]}\\
		&=\dfrac{-7}{12\sqrt{2}}. 
	\end{align*}
	Do đó $  \lim\limits_{x\to 1}\dfrac{\sqrt{x^2+x+2}-\sqrt[3]{7x+1}}{\sqrt{2}(x-1)}=I+J=\dfrac{\sqrt{2}}{12}$.\\
	Suy ra $a=1$, $b=12$,  $c=0$.\\
	Vậy $a+b+c=13$. 
	 }
\end{ex}

% nguồn: Ok_TuanLuan_Toan11_Lan2
\begin{ex}%[1D3V2-5]%[ma:1d8f3bad]
	Giá trị của giới hạn $ \lim\limits_{x\to 0}\dfrac{\left(x^2+{\pi^{21}} \right)\sqrt[7]{1-2x}-{\pi^{21}}}{x}$  là bao nhiêu?
	\loigiai{ 
	 Ta có
	\begin{align*}
		\lim\limits_{x\to 0}\dfrac{\left(x^2+{{\pi }^{21}}  \right)\sqrt[7]{1-2x}-\pi^{21}}{x}&=\lim\limits_{x\to 0}\dfrac{\left(x^2+{{\pi }^{21}} \right)\sqrt[7]{1-2x}-x^2-\pi^{21}+x^2}{x}\\
		&= \lim\limits_{x\to 0}\dfrac{\left( x^2+{{\pi }^{21}} \right)\left( \sqrt[7]{1-2x}-1 \right)}{x}+ \lim\limits_{x\to  0} x\\
		&=\lim\limits_{x\to 0}\dfrac{\left(x^2+\pi^{21}\right)(-2x)}{x\left(\sqrt[7]{1-2x}^6+\sqrt[7]{1-2x}^5+\ldots+\sqrt[7]{1-2x}+1  \right)}\\
		 &=\dfrac{\pi^{21}\cdot (-2)}{1+1+1+1+1+1+1}\\
		&=-\dfrac{2{{\pi  }^{21}}}{7}.
	\end{align*}
	 }
\end{ex}

% nguồn: Ok_TuanLuan_Toan11_Lan2
\begin{ex}%[1D3V2-5]%[ma:6c3f3129]
	Tính giới hạn $L=  \lim\limits_{x\to 1}\dfrac{1-x}{\sqrt{2-x}-1}$.
	\loigiai{ 
	$L=\lim\limits_{x\to 1}\dfrac{1-x}{\sqrt{2-x}-1}=\lim\limits_{x\to 1}\dfrac{(1-x)\left(\sqrt{2-x}+1 \right)}{-x+1}= \lim\limits_{x\to 1}\left(\sqrt{2-x}+1 \right)=2$. 
	 }
\end{ex}

% nguồn: Ok_TuanLuan_Toan11_Lan2
\begin{ex}%[1D3V2-5]%[ma:1d13c788]
	Tính  giới hạn $ \lim\limits_{x\to 3}\,\dfrac{\sqrt{x+1}-\sqrt[3]{x+5}}{x-3}$.
	 \loigiai{
	Ta có 
	\begin{align*}
		 \lim\limits_{x\to 3}\dfrac{\sqrt{x+1}-\sqrt[3]{x+5}}{x-3}&= \lim\limits_{x\to 3}\left( \dfrac{\sqrt{x+1}-2}{x-3}-\dfrac{\sqrt[3]{x+5}-2}{x-3} \right).\\
		&= \lim\limits_{x\to  3}\,\left(\dfrac{x+1-4}{\left( x-3 \right)\left(\sqrt{x+1}+2 \right)}-\dfrac{x+5-8}{(x-3)\left(\sqrt[3]{(x+5)^2}+2\sqrt[3]{x+5}+4 \right)} \right)\\
		&=\lim\limits_{x\to 3}\,\left( \dfrac{1}{\sqrt{x+1}+2}-\dfrac{1}{\sqrt[3]{(x+5)^2}+2\sqrt[3]{x+5}+4} \right)\\ 
		&=\dfrac{1}{4}-\dfrac{1}{12}\\ 
		&=\dfrac{1}{6}. 
	\end{align*}
	 }
\end{ex}

% nguồn: Ok_TuanLuan_Toan11_Lan2
\begin{ex}%[1D3V2-5]%[ma:35f9c4ed]
	Giá trị của  giới hạn $\lim\limits_{x\to +\infty}\left(\sqrt[3]{2x-1}-\sqrt[3]{2x+1}\right)$ là
	 \loigiai{
		Ta  có
		\allowdisplaybreaks 
		\begin{eqnarray*}
			\lim\limits_{x\to +\infty}\left(\sqrt[3]{2x-1}-\sqrt[3]{2x+1}\right)&=&\lim\limits_{x\to +\infty}\dfrac{-2}{\sqrt[3]{\left(2x-1\right)^2}+\sqrt[3]{\left(2x-1\right)\left(2x+1\right)}+\sqrt[3]{\left(2x+1\right)^2}}\\ 
			&=&0. 
		\end{eqnarray*}
	 }
\end{ex}

% nguồn: Ok_TuanLuan_Toan11_Lan2
\begin{ex}%[1D3V2-5]%[ma:9de34648]
	Tính giá trị của giới hạn $ \lim\limits_{x\to 1}\dfrac{\sqrt[3]{x}-1}{\sqrt[3]{4x+4}-2}$. 
	 \loigiai{
	 Ta có
	\begin{align*}
		\lim\limits_ {x\to 1}\dfrac{\sqrt[3]{x}-1}{\sqrt[3]{4x+4}-2}&= \lim\limits_{x\to 1}\dfrac{(x-1)\left(\sqrt[3]{(4x+4)^2}+2\sqrt[3]{4x+4}+4\right)}{(4x+4-8)\left(\sqrt[3]x^2+\sqrt[3]{x}+1  \right)}\\
		&= \lim\limits_{x\to 1}\dfrac{\sqrt[3]{(4x+4)^2}+2\sqrt[3]{4x+4}+4}{4\left( \sqrt[3]x^2+\sqrt[3]{x}+1  \right)}\\
		 &=\dfrac{12}{12}\\
		 &=1.
	\end{align*}
	 }
\end{ex}

% nguồn: Ok_TuanLuan_Toan11_Lan2
\begin{ex}%[1D3V2-5]%[ma:2c9cd078]
	Biết rằng $b>0$, $a+b=5$ và $ \lim\limits_{x\to 0}\,\dfrac{\sqrt[3]{ax+1}-\sqrt{1-bx}}{x}=2$. Tìm  $a$, $b$?
	 \loigiai{
	Ta có 
	\begin{align*}
		\lim\limits_  {x\to 0}\dfrac{\sqrt[3]{ax+1}-\sqrt{1-bx}}{x}&= \lim\limits_{x\to 0}\left( \dfrac{\sqrt[3]{ax+1}-1}{x}+\dfrac{1-\sqrt{1-bx}}{x} \right)\\
		&= \lim\limits_{x\to 0}\left(\dfrac{ax}{x\left(\sqrt[3]{(x+1)^2}+\sqrt[3]{x+1}+1 \right)}+\dfrac{bx}{x\left( 1+\sqrt{1-x} \right)} \right) \\ 
		&= \lim\limits_{x\to 0}\,\left( \dfrac{a}{\left(  \sqrt[3]{(x+1)^2}+\sqrt[3]{x+1}+1 \right)}+\dfrac{b}{\left(1+\sqrt{1-x} \right)}\right)\\
		 &=\dfrac{a}{3}+\dfrac{b}{2}\\
		 &=2.
	\end{align*}
	Vậy ta được $\heva{&a+b=5\\ &\dfrac{a}{3}+\dfrac{b}{2}=2} \Leftrightarrow \heva{&a+b=5  \\&2a+3b=12}\Leftrightarrow \heva{&a=3\\&b=2.}$
	} 
\end{ex}

% nguồn: Ok_TuanLuan_Toan11_Lan2
\begin{ex}%[1D3V2-5]%[ma:fbb099a8]
	Tính giá trị của giới hạn  $ \lim\limits_{x\to 0}\dfrac{2\sqrt{1+x}-\sqrt[3]{8-x}}{x}$.
	\loigiai{ 
	 Ta có
	\begin{align*}
		\lim\limits_{x\to 0}\dfrac{2\sqrt{1+x}-\sqrt[3]{8-x}}{x}&= \lim\limits_{x\to 0}\left(  \dfrac{2\sqrt{1+x}-2}{x}+\dfrac{2-\sqrt[3]{8-x}}{x} \right)\\
		&=  \lim\limits_{x\to 0}\left( \dfrac{2}{\sqrt{x+1}+1}+\dfrac{1}{4+2\sqrt[3]{8-x}+\sqrt[3]{(8-x)^2}} \right)\\
		 &=1+\dfrac{1}{12}\\
		 &=\dfrac{13}{12}.
	\end{align*}
	 }
\end{ex}

% nguồn: Ok_TuanLuan_Toan11_Lan2
\begin{ex}%[1D3V2-5]%[ma:6e8a1ecb]
	 Tính giới hạn $ \lim\limits_{x\to 0}\,\dfrac{\sqrt{4x^2-2x+1}-\sqrt{1-2x}}{x}$.
	\loigiai{ 
	Ta có 
	\begin{align*}
		\lim\limits_{x\to  0}\,\dfrac{\sqrt{4x^2-2x+1}-\sqrt{1-2x}}{x}&= \lim\limits_{x\to 0}\,\dfrac{4x^2}{x\left( \sqrt{4x^2-2x+1}+\sqrt{1-2x} \right)}\\
		&= \lim\limits_{x\to  0}\,\dfrac{4x}{\sqrt{4x^2-2x+1}+\sqrt{1-2x}}\\
		&=0. 
	\end{align*}
	} 
\end{ex}

% nguồn: Ok_TuanLuan_Toan11_Lan2
\begin{ex}%[1D3V2-5]%[ma:1b14051b]
	 Tính giá trị của giới hạn $ \lim\limits_{x\to 1}\,\dfrac{\sqrt{x+3}-2}{x-1}$.
	\loigiai{ 
	Ta có: $ \lim\limits_{x\to 1}\dfrac{\sqrt{x+3}-2}{x-1}= \lim\limits_{x\to  1}\dfrac{x+3-4}{\left( x-1 \right)\left(\sqrt{x+3}+2\right)}=\lim\limits_{x\to 1}\dfrac{1}{\sqrt{x+3}+2}=\dfrac{1}{4}$.
	 }
\end{ex}
